\documentclass[10pt]{article}

% Self-contained NeurIPS-approximating preamble (robust; no external .sty).
\usepackage[letterpaper,margin=1in]{geometry}
\usepackage[T1]{fontenc}
\usepackage[utf8]{inputenc}
\usepackage{lmodern}
\usepackage{amsmath,amssymb,amsfonts,amsthm,mathtools}

% ---- map literal Unicode math symbols to LaTeX (robustness for writer output) ----
\DeclareUnicodeCharacter{03B1}{\ensuremath{\alpha}}
\DeclareUnicodeCharacter{03B2}{\ensuremath{\beta}}
\DeclareUnicodeCharacter{03B3}{\ensuremath{\gamma}}
\DeclareUnicodeCharacter{03B4}{\ensuremath{\delta}}
\DeclareUnicodeCharacter{03B5}{\ensuremath{\varepsilon}}
\DeclareUnicodeCharacter{03B8}{\ensuremath{\theta}}
\DeclareUnicodeCharacter{03BA}{\ensuremath{\kappa}}
\DeclareUnicodeCharacter{03BB}{\ensuremath{\lambda}}
\DeclareUnicodeCharacter{03BC}{\ensuremath{\mu}}
\DeclareUnicodeCharacter{03BD}{\ensuremath{\nu}}
\DeclareUnicodeCharacter{03C0}{\ensuremath{\pi}}
\DeclareUnicodeCharacter{03C1}{\ensuremath{\rho}}
\DeclareUnicodeCharacter{03C3}{\ensuremath{\sigma}}
\DeclareUnicodeCharacter{03C4}{\ensuremath{\tau}}
\DeclareUnicodeCharacter{03C6}{\ensuremath{\varphi}}
\DeclareUnicodeCharacter{03C8}{\ensuremath{\psi}}
\DeclareUnicodeCharacter{03C9}{\ensuremath{\omega}}
\DeclareUnicodeCharacter{0394}{\ensuremath{\Delta}}
\DeclareUnicodeCharacter{0398}{\ensuremath{\Theta}}
\DeclareUnicodeCharacter{039B}{\ensuremath{\Lambda}}
\DeclareUnicodeCharacter{03A0}{\ensuremath{\Pi}}
\DeclareUnicodeCharacter{03A3}{\ensuremath{\Sigma}}
\DeclareUnicodeCharacter{03A6}{\ensuremath{\Phi}}
\DeclareUnicodeCharacter{03A8}{\ensuremath{\Psi}}
\DeclareUnicodeCharacter{03A9}{\ensuremath{\Omega}}
\DeclareUnicodeCharacter{2211}{\ensuremath{\sum}}
\DeclareUnicodeCharacter{220F}{\ensuremath{\prod}}
\DeclareUnicodeCharacter{222B}{\ensuremath{\int}}
\DeclareUnicodeCharacter{221E}{\ensuremath{\infty}}
\DeclareUnicodeCharacter{221D}{\ensuremath{\propto}}
\DeclareUnicodeCharacter{221A}{\ensuremath{\sqrt{}}}
\DeclareUnicodeCharacter{2264}{\ensuremath{\leq}}
\DeclareUnicodeCharacter{2265}{\ensuremath{\geq}}
\DeclareUnicodeCharacter{2260}{\ensuremath{\neq}}
\DeclareUnicodeCharacter{2248}{\ensuremath{\approx}}
\DeclareUnicodeCharacter{2208}{\ensuremath{\in}}
\DeclareUnicodeCharacter{2209}{\ensuremath{\notin}}
\DeclareUnicodeCharacter{2282}{\ensuremath{\subset}}
\DeclareUnicodeCharacter{2286}{\ensuremath{\subseteq}}
\DeclareUnicodeCharacter{2200}{\ensuremath{\forall}}
\DeclareUnicodeCharacter{2203}{\ensuremath{\exists}}
\DeclareUnicodeCharacter{2207}{\ensuremath{\nabla}}
\DeclareUnicodeCharacter{2202}{\ensuremath{\partial}}
\DeclareUnicodeCharacter{2192}{\ensuremath{\rightarrow}}
\DeclareUnicodeCharacter{21A6}{\ensuremath{\mapsto}}
\DeclareUnicodeCharacter{00D7}{\ensuremath{\times}}
\DeclareUnicodeCharacter{2297}{\ensuremath{\otimes}}
\DeclareUnicodeCharacter{2295}{\ensuremath{\oplus}}
\DeclareUnicodeCharacter{2225}{\ensuremath{\|}}
\DeclareUnicodeCharacter{00B7}{\ensuremath{\cdot}}
\DeclareUnicodeCharacter{22C5}{\ensuremath{\cdot}}
\DeclareUnicodeCharacter{00B1}{\ensuremath{\pm}}
\DeclareUnicodeCharacter{211D}{\ensuremath{\mathbb{R}}}
\DeclareUnicodeCharacter{2115}{\ensuremath{\mathbb{N}}}
\DeclareUnicodeCharacter{2124}{\ensuremath{\mathbb{Z}}}
\DeclareUnicodeCharacter{211A}{\ensuremath{\mathbb{Q}}}
\DeclareUnicodeCharacter{2113}{\ensuremath{\ell}}
\DeclareUnicodeCharacter{2605}{\ensuremath{\star}}
\DeclareUnicodeCharacter{22C6}{\ensuremath{\star}}
\DeclareUnicodeCharacter{2227}{\ensuremath{\land}}
\DeclareUnicodeCharacter{2228}{\ensuremath{\lor}}
\DeclareUnicodeCharacter{00AC}{\ensuremath{\neg}}
\DeclareUnicodeCharacter{2308}{\ensuremath{\lceil}}
\DeclareUnicodeCharacter{2309}{\ensuremath{\rceil}}
\DeclareUnicodeCharacter{230A}{\ensuremath{\lfloor}}
\DeclareUnicodeCharacter{230B}{\ensuremath{\rfloor}}
\DeclareUnicodeCharacter{2026}{\ldots}
\DeclareUnicodeCharacter{2212}{\ensuremath{-}}
\DeclareUnicodeCharacter{2032}{\ensuremath{\prime}}
\DeclareUnicodeCharacter{2061}{}
\DeclareUnicodeCharacter{2062}{}
\usepackage{natbib}
\usepackage{booktabs}
\usepackage{microtype}
\usepackage{xcolor}
\usepackage{graphicx}
\usepackage{caption}
\usepackage{enumitem}
\usepackage{listings}
\usepackage[colorlinks=true,linkcolor=blue,citecolor=blue,urlcolor=blue]{hyperref}
\usepackage{url}
\usepackage[capitalize]{cleveref}
\usepackage{setspace}

\bibliographystyle{plainnat}
\setlength{\parskip}{2pt}

% ---- theorem environments ----
\theoremstyle{plain}
\newtheorem{theorem}{Theorem}
\newtheorem{lemma}{Lemma}
\newtheorem{corollary}{Corollary}
\newtheorem{proposition}{Proposition}
\theoremstyle{definition}
\newtheorem{definition}{Definition}
\newtheorem{assumption}{Assumption}
\theoremstyle{remark}
\newtheorem{remark}{Remark}

% cleveref names for custom theorem-like environments
\crefname{assumption}{Assumption}{Assumptions}
\Crefname{assumption}{Assumption}{Assumptions}
\crefname{definition}{Definition}{Definitions}
\Crefname{definition}{Definition}{Definitions}
\crefname{proposition}{Proposition}{Propositions}
\Crefname{proposition}{Proposition}{Propositions}
\crefname{remark}{Remark}{Remarks}
\Crefname{remark}{Remark}{Remarks}

% ---- shared notation macros (\ensuremath so they work in text OR math mode) ----
\newcommand{\cZ}{\ensuremath{\mathcal{Z}}}
\newcommand{\cZA}{\ensuremath{\mathcal{Z}_{\mathrm{A}}}}
\newcommand{\cZB}{\ensuremath{\mathcal{Z}_{\mathrm{B}}}}
\newcommand{\qo}{\ensuremath{q_{0}}}
\newcommand{\qz}{\ensuremath{q}}
\newcommand{\qs}{\ensuremath{q^{\star}}}
\newcommand{\Rwd}{\ensuremath{R}}
\newcommand{\Zpart}{\ensuremath{Z}}
\newcommand{\gain}{\ensuremath{\mathcal{G}}}
\newcommand{\spread}{\ensuremath{\mathrm{spread}}}
\newcommand{\Exp}{\ensuremath{\mathbb{E}}}
\newcommand{\KL}[2]{\ensuremath{\mathrm{KL}\!\left(#1 \,\|\, #2\right)}}
\newcommand{\Fobj}[1]{\ensuremath{F\!\left(#1\right)}}
\newcommand{\R}{\ensuremath{\mathbb{R}}}
% safety aliases sometimes used by writers
\providecommand{\E}{\ensuremath{\mathbb{E}}}
\providecommand{\reals}{\ensuremath{\mathbb{R}}}
\providecommand{\bbR}{\ensuremath{\mathbb{R}}}
\providecommand{\bbN}{\ensuremath{\mathbb{N}}}
\providecommand{\bbZ}{\ensuremath{\mathbb{Z}}}
\providecommand{\bbQ}{\ensuremath{\mathbb{Q}}}
\providecommand{\bbE}{\ensuremath{\mathbb{E}}}
\providecommand{\bbP}{\ensuremath{\mathbb{P}}}
\providecommand{\Tr}{\ensuremath{\mathrm{Tr}}}
\providecommand{\diag}{\ensuremath{\mathrm{diag}}}
\providecommand{\sgn}{\ensuremath{\mathrm{sgn}}}
\providecommand{\eqdef}{\ensuremath{\coloneqq}}
\DeclareMathOperator*{\argmax}{arg\,max}
\DeclareMathOperator*{\argmin}{arg\,min}

\lstset{basicstyle=\ttfamily\footnotesize,breaklines=true,columns=fullflexible,
  keepspaces=true,showstringspaces=false,frame=single,framesep=3pt}

\title{Training-Free Reinforcement Learning Has a Vanishing Optimization
Space on Single Models but Recovers in Context-Isolated Agent Systems:\\[2pt]
\large Theory and a Self-Evolving Omni Speech Agent}

\author{Exploring-L4-Intelligence Project}
\date{}

\begin{document}
\maketitle


% ===== section 01-intro (writer wordcount 1180) =====
\begin{abstract}
Modern omni speech LLMs absorb broad cross-modal, multi-granularity knowledge during pretraining, yet at inference time the levers a practitioner controls---instructions, demonstrations, decoding---induce a strikingly small, frequently \emph{flat} space of achievable behaviors. We make this precise. Casting training-free reinforcement learning as Gibbs tilting of a frozen base policy, $\qs \propto \qo\exp(\Rwd/\beta)$, we prove an \emph{optimization-space adequacy} (OSA) trichotomy, machine-checked in Lean~4. On a single model the realizable reward spread $\spread$ collapses, so the KL-regularized gain $\gain$ vanishes: inference-time RL is asymptotically inert (OSA-1). Lifting the same frozen model into a \emph{context-isolated agent system} renders the action space separable across sub-tasks, making the gain additive and growing in the number of isolated contexts (OSA-2)---at the cost of rollout instability that only algorithm-level credit assignment under a $\beta$-KL trust region tames (OSA-3). We instantiate the theory as a self-evolving omni \emph{speech} agent pairing a frozen vector-omni memory, a generative-omni policy, and a verifiable-reward acceptance gate. We are candid that the agent \emph{mechanism} is not novel; the contribution is the domain transfer, the two-omni pairing, verifiable speech rewards, and the optimization-space theory that explains why such systems must replace single-model search. We report preliminary feasibility results and a research plan, including a proposed cross-session paralinguistic memory benchmark.
\end{abstract}

\section{Introduction}
\label{sec:intro}

A modern omni or speech-capable multimodal LLM is, by the time it is released, an extraordinary repository of latent capability. From large-scale unsupervised audio, parallel speech-text corpora, and instruction mixtures, models such as Qwen2-Audio, Qwen2.5-Omni, Qwen3-Omni, and SALMONN \citep{qwen2audio,qwen25omni,qwen3omni,salmonn2023} acquire transcription, translation, speaker identity, paralinguistic affect, and dialogue competence simultaneously. The animating question of this program is deliberately conservative: \emph{how far can one go without touching a single weight?} We treat the pretrained model as frozen and ask what reward-guided optimization, applied purely at inference time, can elicit. This is the regime of \emph{training-free} reinforcement learning---best-of-$N$ selection, reward-guided decoding, reranking, task-conditioning, and representation read-out---that changes no parameters and no architecture \citep{beirami2024bon,mudgal2024controlled,zuo2025ttrl}.

\paragraph{The action space is small, and often flat.}
The premise that motivates everything below is empirical and, we will argue, structural. Although the frozen model \emph{contains} broad capability, the degrees of freedom a practitioner actually controls at inference---the prompt, the demonstrations, the sampling temperature---move the model through a narrow envelope of behaviors. In-context demonstrations frequently fix \emph{format} rather than \emph{accuracy}: their labels can be permuted or corrupted with little effect on the answer distribution \citep{min2022rethinking,alice2026}, and for speech in particular, models tend to ``read'' a transcribable surface rather than ``listen'' to the acoustic substrate, so paralinguistic conditioning barely perturbs the output \citep{voxparadox2026,hsu2023iclspeech}. Our own in-house measurements corroborate this: across task-conditioning variants on a frozen omni backbone, the conditional output distribution is close to flat in the reward dimension---the candidate set generated by any single-model search collapses onto a thin band of near-identical rewards. We name this phenomenon \emph{flat conditioning}, and it is the central obstacle that the theory in this paper is built to characterize.

\paragraph{Training-free RL as Gibbs tilting.}
To reason about the obstacle quantitatively we adopt the now-standard variational view of inference-time alignment. Optimizing the KL-regularized objective
\[
\Fobj{\qz} \;=\; \Exp_{z\sim \qz}[\Rwd(z)] \;-\; \beta\,\KL{\qz}{\qo}
\]
over search distributions $\qz$ on the inference-time action space $\cZ$ has the closed-form Gibbs optimum $\qs(z)\propto \qo(z)\exp(\Rwd(z)/\beta)$, with partition function $\Zpart$ \citep{korbak2022kl,beirami2024bon}. Every practical training-free method---best-of-$N$ and its soft and asymptotic refinements \citep{verdun2025softbon,yang2024asymptotics,gui2024bonbon}, self-consistency \citep{wang2023selfconsistency}, MBR decoding \citep{ichihara2025mbr,mbrasr2025}, and controlled decoding \citep{mudgal2024controlled}---is an estimator of this tilt over some realization of $\cZ$. The attainable benefit of \emph{any} such estimator is therefore bounded by how far $\qs$ can depart from $\qo$, which in turn is governed by the spread of $\Rwd$ that the base distribution $\qo$ actually places mass on. Flat conditioning is precisely the statement that this spread is small.

\paragraph{Central thesis: optimization-space adequacy.}
Our contribution is to turn this observation into a sharp, verifiable trichotomy that we call \emph{optimization-space adequacy} (OSA). \textbf{OSA-1 (single-model collapse).} When $\Rwd$ is concentrated under $\qo$---reward spread $\spread\to 0$---the Gibbs tilt is a vanishing perturbation and the optimization gain obeys $\gain = \Fobj{\qs}-\Fobj{\qo}\le \tfrac{1}{2\beta}\spread^{2}\to 0$. Inference-time RL on a single frozen model is asymptotically inert: more samples, more clever reranking, and more elaborate search cannot manufacture headroom that the conditional distribution does not contain. This mirrors, from the optimization side, the reward-model over-optimization scaling laws observed empirically \citep{gao2023overopt}, and we prove it as a clean inequality, machine-checked in Lean~4. \textbf{OSA-2 (agentic recovery).} The escape is structural, not numerical. Decomposing a task across a \emph{context-isolated agent system}---where each isolated context exposes an independent operator on $\cZ$, whether the embedding/Operator-A view $\cZA$ or the generative/Operator-B view $\cZB$---makes the joint action space a product, and the attainable gain becomes \emph{additive} across isolated contexts rather than capped at the single-model ceiling. Isolation is what restores spread: separate contexts do not average one another's rewards toward the flat band. \textbf{OSA-3 (the stability tax).} This recovery is not free. A product action space optimized by rollout incurs variance and feedback-loop pathologies---context collapse, error propagation in append-only memories, and ``memory that hurts'' \citep{ace2025,expfollowing2025,remembering2026}---so the additive gain is realizable only under algorithm-level credit assignment \citep{jitrl2026,ahmadian2024rloo} and an explicit $\beta$-KL trust region that keeps each context's search distribution close to its own base. OSA-1, OSA-2, and OSA-3 together say: \emph{do not search harder on one model; isolate, decompose, and stabilize.}

\paragraph{A self-evolving omni speech agent.}
We instantiate the theory in the most demanding modality. The system pairs (i) a \emph{frozen vector-omni memory}---a bi-encoder omni-embedding model \citep{omniembed2025,omniembedaudio2026} that stores task-conditioned acoustic views as durable, retrievable context; (ii) a \emph{generative-omni policy} that acts under retrieved context; and (iii) a \emph{verifiable-reward acceptance gate} that admits a candidate skill or memory only when an automatically checkable speech reward (WER for ASR, verifiable matches for translation and intent, and learned speech reward models for paralinguistics) certifies improvement \citep{tulu3,rlvrincentive2025,gsrm2026,paras2s2025}. The two omni models---one as content-addressable memory, one as actor---are the context-isolation substrate that OSA-2 requires, and the acceptance gate is the trust region OSA-3 demands. To our knowledge no speech-native, no-gradient, self-evolving agent of this form exists; the omni-embedding-as-memory pairing is, at present, an open moat.

\paragraph{What is and is not novel.}
We are explicit and honest about positioning. The \emph{agent mechanism}---reflective self-improvement, skill libraries, memory-augmented policies---is not novel; the general-agent frontier is crowded with strong, closely related designs \citep{reflexion2023,voyager2023,jitrl2026,mem0_2025,memgpt2023}. Our novelty lies elsewhere: in the \emph{domain transfer} of these ideas to frozen omni speech models, in the \emph{two-omni pairing} of a vector memory with a generative policy, in \emph{verifiable speech rewards} as the acceptance criterion, and---most importantly---in the \emph{optimization-space theory} that explains \emph{why} a single-model training-free approach must give way to an agentic one. The theory is the load-bearing contribution; the system is its existence proof.

\paragraph{Contributions.}
We make five.
\textbf{(C1) Theory.} We formalize optimization-space adequacy and prove the OSA-1/OSA-2/OSA-3 trichotomy---single-model gain collapse, additive agentic recovery, and the stability tax---with the core inequalities machine-checked in Lean~4.
\textbf{(C2) Convergence analysis.} We connect OSA-3 to a survey of agentic instability (context collapse, append-only error propagation, memory-induced regression) and give conditions, under $\beta$-KL trust regions and credit assignment, for the agent loop to converge rather than plateau or diverge \citep{ace2025,expfollowing2025,remembering2026}.
\textbf{(C3) Feasibility, moat, and preliminary results.} We argue feasibility, document the open moat for speech-native no-gradient self-evolution, and report preliminary measurements of flat conditioning and of early agentic recovery.
\textbf{(C4) System design.} We specify the self-evolving omni speech agent---frozen vector-omni memory, generative-omni policy, verifiable-reward gate---as a concrete, reproducible architecture.
\textbf{(C5) Research plan and benchmark.} We lay out a staged research plan and propose a \emph{cross-session paralinguistic memory} benchmark for speech agents, addressing a gap in current memory and multi-turn audio evaluations \citep{longmemeval2024,locomo2024,audiomc2025}.

\paragraph{Organization.}
\Cref{sec:intro} has framed flat conditioning and the OSA thesis. The remainder of the paper develops the formal setting and the Gibbs-tilt objective; states and proves OSA-1 (single-model collapse), OSA-2 (agentic recovery under context isolation), and OSA-3 (the stability tax) with the Lean-verified inequalities; analyzes convergence of the agent loop under trust-region credit assignment; presents the self-evolving omni speech agent and its verifiable-reward gate together with feasibility and preliminary results; introduces the proposed cross-session paralinguistic memory benchmark; and closes with related work, limitations, and the staged research plan.


% ===== section 02-prelim (writer wordcount 1180) =====
\section{Preliminaries and Problem Formulation}
\label{sec:prelim}

We study \emph{training-free} reinforcement learning (RL): a regime in which a frozen multimodal speech/omni model is steered at inference time by a reward signal, with no gradient ever propagated into its parameters and no architectural surgery. This section fixes notation, states the invariant that distinguishes training-free RL from conventional alignment, isolates the two model classes and their associated inference-time operators, lifts the formulation to an agentic action space with context isolation, and specifies the reward class we admit. The objects defined here are reused verbatim in the Theory section.

\subsection{Training-free RL as constrained inference-time reward maximization}
\label{subsec:tfrl}

Let $g_\theta$ denote a pretrained speech/omni model with \emph{frozen} parameters $\theta$, and let $x$ be an input (an utterance, a prompt, or a paired context). The inference-time degrees of freedom---the only quantities we are permitted to vary---are collected into an \emph{action} $z \in \cZ$, where $\cZ$ is the model- and task-dependent action space made precise in \Cref{subsec:classes,subsec:agent}. Executing an action produces an output $g_\theta(x,z)$ that is scored by a bounded scalar \emph{reward} $\Rwd\colon \cZ \to [0,1]$ (we write $\Rwd(z)$ for $\Rwd(g_\theta(x,z))$ when $x$ is fixed). A pretrained \emph{base} (or prior) distribution $\qo \in \Delta(\cZ)$ over actions is induced by the frozen model and the default inference procedure---for a generative model, the temperature-scaled autoregressive sampler; for an embedding model, the canonical encoding policy. Training-free RL searches over distributions $\qz \in \Delta(\cZ)$ for one that raises expected reward while remaining close to the prior.

\begin{definition}[Training-free RL objective]
\label{def:objective}
Fix $x$, a base distribution $\qo$, a reward $\Rwd$, and a temperature $\beta > 0$. The KL-regularized inference-time objective is
\begin{equation}
\label{eq:Fobj}
\Fobj{\qz} \;=\; \Exp_{z \sim \qz}\!\big[\Rwd(z)\big] \;-\; \beta\,\KL{\qz}{\qo},
\end{equation}
and the training-free RL problem is $\qs \;=\; \arg\max_{\qz \in \Delta(\cZ)} \Fobj{\qz}$.
\end{definition}

The objective \eqref{eq:Fobj} is strictly concave in $\qz$ over the simplex, and its maximizer is the classical Gibbs/exponential tilt of the prior \citep{korbak2022kl}.

\begin{proposition}[Gibbs optimum]
\label{prop:gibbs}
The unique maximizer of \eqref{eq:Fobj} is
\begin{equation}
\label{eq:gibbs}
\qs(z) \;=\; \frac{1}{\Zpart}\,\qo(z)\,\exp\!\Big(\tfrac{1}{\beta}\,\Rwd(z)\Big),
\qquad
\Zpart \;=\; \Exp_{z \sim \qo}\!\Big[\exp\!\big(\tfrac{1}{\beta}\Rwd(z)\big)\Big],
\end{equation}
with optimal value $\Fobj{\qs} = \beta \log \Zpart$.
\end{proposition}

\begin{proof}[Proof sketch]
Adding a Lagrange multiplier for the normalization constraint and differentiating the concave functional \eqref{eq:Fobj} with respect to $\qz(z)$ yields $\Rwd(z) - \beta(\log \qz(z) - \log \qo(z)) - \beta = \text{const}$, i.e. $\log \qz(z) = \log \qo(z) + \Rwd(z)/\beta + \text{const}$; exponentiating and normalizing gives \eqref{eq:gibbs}. Substituting back collapses the objective to $\beta\log\Zpart$.
\end{proof}

Equation~\eqref{eq:gibbs} is the formal anchor of the paper: every concrete inference-time procedure we analyze---best-of-$N$, soft best-of-$N$, minimum-Bayes-risk decoding, reward-guided decoding, and search over embedding configurations---is read as an \emph{estimator of, or a sampler from, the tilted target} $\qs$, and is therefore comparable through the lens of how well it approximates \eqref{eq:gibbs} at a given query budget. The temperature $\beta$ interpolates between the prior ($\beta\to\infty$, $\qs\to\qo$) and greedy reward maximization ($\beta\to 0^+$, $\qs\to$ the $\arg\max_z \Rwd$ point mass), recovering best-of-$N$ alignment in the low-temperature limit \citep{beirami2024bon,verdun2025softbon,khalaf2025hedgetune}.

\begin{definition}[No-weight-update invariant]
\label{def:invariant}
A procedure is \emph{training-free} if, for every input and every query budget, the model parameters $\theta$ and the model's computational graph are held fixed; all adaptation is realized through the choice of $\qz$ over the inference-time action space $\cZ$. Formally, the map $\theta \mapsto \theta$ is the identity throughout, and only $\qz$ is optimized.
\end{definition}

\Cref{def:invariant} is the operational contract of this work. It excludes parameter-efficient fine-tuning and any backpropagation through $g_\theta$; it admits only \emph{selection, reweighting, search, and prompt/context construction} over $\cZ$. Because $\qs$ depends on $\theta$ only through $\qo$ and on the task only through $\Rwd$, the frozen model's pretrained competence is \emph{activated}, never overwritten---the central premise the paper develops empirically and theoretically.

\subsection{Two model classes and their inference-time operators}
\label{subsec:classes}

The action space $\cZ$ factorizes by model class. We distinguish a \emph{vector} (embedding) class and a \emph{generative} (thinker--talker) class, with corresponding operators $\cZA$ and $\cZB$, $\cZ = \cZA \cup \cZB$.

\paragraph{Vector / embedding omni (Operator A).}
A contrastive bi-encoder omni model maps a (conditioned) input to a single $\ell_2$-normalized vector, typically by masked-mean pooling over a chosen hidden layer followed by a projection \citep{omniembed2025,clap2022,nvembed2024}. Here $g_\theta$ emits a representation rather than a token sequence, the reward is a representation-quality functional (retrieval hit, probe accuracy, alignment/uniformity surrogate), and the prior $\qo$ is the default encoding policy.

\begin{definition}[Operator A: embedding search space $\cZA$]
\label{def:opA}
For the vector class, an action $z_{\mathrm A} = (c, \ell, \pi, P, S) \in \cZA$ specifies (i) the \emph{conditioning} $c$ (instruction/prompt prepended to the audio), (ii) the \emph{layer} $\ell$ whose hidden states are read, (iii) the \emph{pooling} rule $\pi$ (masked mean, attentive statistics, last-token), (iv) the \emph{projection} $P$ applied to the pooled vector, and (v) the \emph{selection} $S$ over candidate views/segments. The output embedding is $u(z_{\mathrm A}) = P\big(\pi_\ell(g_\theta(x,c))\big)/\lVert\cdot\rVert_2$. Operator A is the search over $\cZA$ that maximizes a verifiable representation reward.
\end{definition}

\paragraph{Generative thinker--talker omni (Operator B).}
An autoregressive omni model defines a distribution over token sequences $y$ conditioned on $(x,c)$; the default sampler is the prior $\qo$ \citep{qwen3omni}. Inference-time alignment reshapes $\qo$ toward $\qs$ either by drawing $N$ candidates and selecting/aggregating (best-of-$N$, minimum-Bayes-risk) or by tilting the per-step transition with a value estimate (reward-guided/controlled decoding) \citep{mudgal2024controlled,beirami2024bon,verdun2025softbon}.

\begin{definition}[Operator B: generative decoding space $\cZB$]
\label{def:opB}
For the generative class, an action $z_{\mathrm B} \in \cZB$ specifies a decoding policy over output sequences: a sampling temperature, a candidate budget $N$, and an aggregation rule $A \in \{\text{best-of-}N,\ \text{MBR},\ \text{reward-guided}\}$. Best-of-$N$ returns $\arg\max_{i\le N}\Rwd(y_i)$ for $y_i\sim\qo$; MBR returns $\arg\max_i \Exp_{y'\sim\qo}[u(y_i,y')]$ for a utility $u$ \citep{ichihara2025mbr,mbrasr2025}; reward-guided decoding tilts each transition by $\exp(V(\cdot)/\beta)$ to approximate \eqref{eq:gibbs} token-by-token \citep{mudgal2024controlled}.
\end{definition}

Operators A and B are not interchangeable: A reweights a \emph{geometric} object (a point on the unit sphere) and contributes no entropy to the output, whereas B reweights a \emph{combinatorial} object (a sequence) drawn from a high-entropy prior. This asymmetry---a single-sample deterministic readout versus an $N$-sample stochastic search---governs the budget--quality trade-offs analyzed later and motivates treating the two as distinct instantiations of the same tilt \eqref{eq:gibbs}.

\subsection{The agent action space and context isolation}
\label{subsec:agent}

When the frozen model is embedded in an agent, the action is no longer a single decode but a \emph{pipeline}: the agent first \emph{recalls} memory, then \emph{invokes} a skill, then \emph{constructs} context, then \emph{decodes}. We make this explicit.

\begin{definition}[Agentic action]
\label{def:agentaction}
An agentic action is a tuple $z = (m, s, c, d)$ where $m$ selects which memory item(s) to recall, $s$ selects which skill/tool to invoke, $c$ builds the context window from $(x,m,s)$, and $d \in \cZA \cup \cZB$ is the terminal embedding or decoding action of \Cref{def:opA,def:opB}. The base distribution $\qo$ factorizes as $\qo(z) = \qo(m)\,\qo(s\mid m)\,\qo(c\mid m,s)\,\qo(d\mid c)$.
\end{definition}

Many agentic systems decompose a task across $K$ sub-agents that operate on disjoint slices of the action. We formalize the structure that makes such a decomposition tractable.

\begin{definition}[Context isolation / separable structure]
\label{def:isolation}
The action space exhibits \emph{context isolation} across $K$ sub-agents if it factorizes as a product $\cZ = \prod_{k=1}^{K}\cZ^{(k)}$ with $z=(z^{(1)},\dots,z^{(K)})$, the base distribution is independent across blocks, $\qo(z)=\prod_{k}\qo^{(k)}(z^{(k)})$, and the reward is \emph{additively separable}, $\Rwd(z)=\sum_{k=1}^{K}\Rwd^{(k)}(z^{(k)})$.
\end{definition}

\begin{remark}
\label{rem:factorize}
Under \Cref{def:isolation}, the Gibbs optimum \eqref{eq:gibbs} factorizes blockwise, $\qs(z)=\prod_{k}{\qs}^{(k)}(z^{(k)})$ with ${\qs}^{(k)}\propto\qo^{(k)}\exp(\Rwd^{(k)}/\beta)$ and $\Zpart=\prod_k \Zpart^{(k)}$. Consequently the search decouples: each sub-agent may sample/select within its own block without cross-block interference, and the partition function---hence the optimization gain---is the product (equivalently, sum in log-space) of per-block contributions. This separability is the formal object the second operator-level result (OSA-2) requires, and it is precisely what \emph{context contamination} between sub-agents destroys.
\end{remark}

\subsection{Verifiable versus model-judged rewards}
\label{subsec:reward}

The validity of every guarantee in \Cref{subsec:tfrl} rests on $\Rwd$ being trustworthy. We admit only \emph{verifiable} rewards.

\begin{definition}[Verifiable reward]
\label{def:verifiable}
A reward $\Rwd$ is \emph{verifiable} if it is a deterministic, label-derived function of the output and a ground-truth annotation $y^\star$, computable without invoking a learned scorer. Canonical instances are (i) negative word error rate $1-\mathrm{WER}(y,y^\star)$ for ASR/ST, (ii) exact-match/accuracy for classification and intent, (iii) retrieval hit$@k$ for embedding queries, and (iv) probe accuracy for a fixed linear probe. A reward is \emph{model-judged} if it is the output of a separate generative or preference model evaluating $y$.
\end{definition}

\begin{assumption}[Verifiability]
\label{ass:verifiable}
Throughout, $\Rwd$ is verifiable in the sense of \Cref{def:verifiable}: bounded in $[0,1]$, deterministic given $y^\star$, and independent of $g_\theta$.
\end{assumption}

We adopt \Cref{ass:verifiable} for two reasons. First, model-judged rewards are susceptible to \emph{reward hacking}---the optimizer exploits gaps between the proxy and the true objective, so that increasing $\Exp_{\qs}[\Rwd]$ no longer increases task quality \citep{skalse2022rewardhacking,gao2023overopt}. Because the tilt \eqref{eq:gibbs} sharpens precisely along the direction of $\Rwd$, any reward misspecification is amplified rather than averaged out as $\beta\to 0$. Second, learned judges carry systematic artifacts---position bias and self-preference---that contaminate the ranking used by Operator B \citep{positionbias2024,selfpreference2024}; verifiable rewards are invariant to these. Restricting to label-derived $\Rwd$ is therefore not a convenience but a precondition for the no-weight-update invariant of \Cref{def:invariant} to translate into genuine, contamination-free activation of pretrained competence \citep{tulu3}. With \Cref{def:objective,def:invariant,def:opA,def:opB,def:isolation,def:verifiable} and \Cref{ass:verifiable} in place, the notation required by the Theory section is complete.


% ===== section 03-related-a (writer wordcount 1187) =====
\section{Related Work}
\label{sec:related}

Our contribution sits at the confluence of three literatures that have, until recently, evolved in relative isolation: the theory of inference-time reward optimization as exponential tilting, the empirical study of in-context learning (ICL) in audio language models, and the geometric characterization of what frozen speech encoders represent. We review each in turn, emphasizing not a catalogue of methods but the structural tensions that motivate our framework: that training-free RL is a search over a fixed action space \cZ whose efficacy is bounded by the support of the base distribution \qo, and that this boundedness manifests differently for the embedding operator \cZA than for the generative operator \cZB.

\subsection{Inference-Time / Training-Free RL and the Tilting View}
\label{sec:rw-tilting}

The unifying abstraction behind training-free RL is that a reward-guided search distribution \qz is, in its optimal form, an exponential tilt of the base policy: the KL-regularized objective \(\Fobj{q}=\Exp_{q}[\Rwd]-\beta^{-1}\KL{q}{\qo}\) is maximized by the Gibbs measure \(\qs(z)\propto\qo(z)\exp(\beta\Rwd(z))\), with partition function \Zpart \citep{korbak2022kl}. This identity recasts a spectrum of seemingly disparate inference-time procedures as approximations to the same tilt. Best-of-$N$ (BoN) sampling is the most studied instance: \citet{beirami2024bon} establish that the KL divergence between the BoN distribution and \qo grows as \(\log N - (N{-}1)/N\), giving a closed-form exchange rate between compute and the distributional shift purchased. \citet{verdun2025softbon} sharpen this with a soft, temperature-controlled BoN whose suboptimality decays as \(O(1/N)\), while \citet{aminian2025smoothing} recast BoN through a smoothing lens, bounding its regret against the Gibbs optimum and exposing when additional samples cease to help. The crucial—and for our purposes load-bearing—observation across this line is that \emph{every} tilt reweights existing samples: the support of \qz is contained in the support of \qo. No amount of $N$ creates mass where \qo assigns none. Minimum Bayes risk (MBR) decoding \citep{ichihara2025mbr}, including its specialization to ASR \citep{mbrasr2025}, and controlled decoding \citep{mudgal2024controlled} inherit the same ceiling: they are consensus or value-guided reweightings, powerful when the correct hypothesis is already sampled with non-negligible probability and inert otherwise.

A more recent strand attempts to extract reward signal without any labels at all. TTRL \citep{zuo2025ttrl} and TPO \citep{li2025tpo} use self-consistency or majority vote as a surrogate reward at test time, and JitRL \citep{jitrl2026} performs gradient-free credit assignment over a frozen policy. These methods are attractive precisely because they touch no weights, but they sharpen rather than dissolve the support constraint: a self-generated reward can only amplify modes the model already prefers, which is why majority-vote rewards can entrench confident errors. The over-optimization literature quantifies the resulting hazard. \citet{gao2023overopt} show that reward and true objective diverge as KL grows, following a predictable scaling law; \citet{khalaf2025hedgetune} characterize the optimal effective $N^\ast$ (the Best-of-Poisson regime) beyond which tilting harms. The over-optimization risk is governed by the reward spread \spread and the realized gain \gain: when the proxy reward is mis-specified, larger \gain trades calibrated base knowledge for spurious peaks.

These dynamics intersect the ongoing RLVR boundary debate, which directly concerns whether inference-time or verifiable-reward optimization can expand a model's capabilities or merely re-allocate them. \citet{rlvrincentive2025} argue that RLVR predominantly sharpens the base distribution rather than adding skills, a position echoed by the finding that even spurious or random rewards can improve benchmark scores by exploiting format and prior structure \citep{spuriousrewards2025}. The countervailing evidence that RLVR can induce genuinely new skills \citep{newskills2026} does not contradict the tilting view so much as locate its exception: new behavior emerges only when the action space \cZ is rich enough—e.g., multi-step generative composition—that the relevant trajectory had non-zero base mass that single-shot sampling rarely surfaced. This distinction—reweighting versus genuine expansion—is exactly the axis along which we separate \cZA from \cZB, and the tilting formalism makes the separation precise rather than rhetorical.

\subsection{In-Context Learning in Audio LLMs and the Two Model Classes}
\label{sec:rw-icl}

If training-free RL reweights the base distribution, the natural question is whether in-context demonstrations can \emph{move} it. The text-domain ICL literature offers a sobering prior. \citet{min2022rethinking} demonstrate that demonstrations are largely label-insensitive: corrupting ground-truth labels barely degrades performance, implying that demonstrations specify the task and output format rather than teaching the input-output mapping. \citet{alice2026} reach a congruent conclusion for instruction-following demonstrations, which fix format but not accuracy. ICL is moreover scale-gated—its more sophisticated, override-the-prior behaviors emerge only in larger models \citep{wei2023larger}—a property that carries into audio, where few-shot competence appears as an emergent phenomenon in larger audio LLMs \citep{mimoaudio2025}. The function- and task-vector view \citep{todd2024functionvectors,taskvectors2025} explains the mechanism: demonstrations compress to a latent task direction that conditions decoding without altering the underlying perceptual representation. For a generative operator \cZB, this means demonstrations can re-route \emph{how} content is verbalized, but they do not re-open perceptual channels the encoder has discarded.

This perceptual ceiling is starkly visible in audio. A growing body of evidence shows that audio LLMs frequently ``read rather than listen'': \citet{voxparadox2026} document that models answer paralinguistic questions from textual priors and transcript content rather than the acoustic signal, and \citet{listenortranscribe2025} find that performance often depends on whether the model transcribes versus genuinely attends to non-lexical cues. In-context demonstrations do little to repair this when the requisite acoustic information has already been pooled away upstream—consistent with the broader finding that few-shot prompting fixes format, not the missing percept. SER is a partial and instructive exception: \citet{sericl2025} report that few-shot demonstrations \emph{do} improve speech emotion recognition, plausibly because emotion is one of the paralinguistic factors a generative omni model retains with enough base probability for tilting and ICL to surface. The pattern—ICL helping for retained factors, failing for discarded ones—mirrors the support constraint of \Cref{sec:rw-tilting} and foreshadows our operator-dependent treatment.

The contrast with embedding models is decisive for our two-class taxonomy. Text embedders do not acquire ICL for free: \citet{bgeenicl2024} show that exploiting in-context examples in an embedder requires dedicated training, and instruction-conditioned embedders such as INSTRUCTOR \citep{su2023instructor} must be explicitly trained to follow task instructions at encode time. An off-the-shelf embedding operator \cZA therefore exposes \emph{no} prompt-conditioned degree of freedom analogous to \cZB's in-context steering: its action space is the choice of pooling, layer, and query, not a re-conditioning of the forward pass. This asymmetry—\cZB admits demonstration-driven tilts, \cZA does not—is, in our framework, not an implementation detail but a defining structural difference between the two model classes, and it dictates which training-free interventions are even available for each.

\subsection{Capability Map, Embedding Geometry, and Probing}
\label{sec:rw-geometry}

Whether a frozen model \emph{can} be tilted toward a task ultimately depends on whether the relevant information survives in its representations—a question the probing and embedding-geometry literatures answer with unusual clarity. Contrastive and alignment objectives are now understood to actively discard nuisance variation: \citet{xiao2020whatnot} show that contrastive learning suppresses factors not aligned with the training augmentations, and \citet{wang2020alignuniform} formalize the alignment–uniformity trade-off that drives semantically similar inputs together while spreading the hypersphere uniformly. For a content-oriented embedder, speaker identity, emotion, and channel are precisely the nuisances such objectives are designed to remove. This is the geometric root of why an embedding operator \cZA, however cleverly searched, cannot recover a factor its training objective erased: the base distribution \qo over \cZA simply has no mass on the discarded direction.

Pooling compounds the loss. Mean- or last-token pooling imposes an information bottleneck that NV-Embed's learned latent-attention pooling was introduced to mitigate \citep{nvembed2024}, and the speaker-verification community has long known that attentive statistics pooling \citep{attnstatspool2018} and ECAPA-TDNN's multi-scale aggregation \citep{ecapatdnn2020} are necessary precisely because naive temporal averaging destroys speaker-discriminative structure. The implication is that the \emph{same} encoder can be near-blind or highly informative for a factor depending solely on the read-out, which is exactly the degree of freedom a training-free search over \cZA should exploit. Layer-wise probing makes this concrete: \citet{pasad2021layerwise} chart an autoencoder-like trajectory across wav2vec~2.0 layers in which phonetic and word-level information peaks at intermediate depths, and \citet{spkattrprobing2025} show speaker attributes are distributed non-uniformly across layers. Capability is thus not a scalar property of a model but a function of (layer, pooling, query)—the very coordinates of our action space \cZA.

Crucially, evidence that paralinguistic content \emph{persists} even when unused is now substantial. \citet{frozenllmparaling2024} demonstrate that a frozen LLM, given speech features, perceives paralinguistic attributes it never verbalizes; \citet{emotionneurons2026} localize emotion-selective units inside such models. These findings establish the optimistic precondition for our flagship study: the information needed to disentangle content, speaker, emotion, and intent is present in the frozen omni model's hidden states, even when the default forward pass—or a generative decode \cZB—suppresses it. The task is therefore one of \emph{activation} via search over read-outs, not of teaching. The SUPERB and Dynamic-SUPERB taxonomies \citep{superb2021,dynsuperb2023} supply the benchmark scaffolding for evaluating such activation across content, speaker, semantic, and paralinguistic axes, and we adopt their factorization to structure our capability map. What this literature lacks—and what our tilting formalism supplies—is a principled account of \emph{which} of these latent capabilities are reachable by reweighting a fixed \qo and which demand a different action space altogether.

% ===== section 04-related-b (writer wordcount 963) =====
\subsection{Agent Memory}
\label{sec:rw-memory}

A memory module is the substrate on which a training-free agent accumulates experience without touching weights, and the literature has moved from flat buffers to increasingly structured stores. Early designs append raw interaction logs and retrieve by recency or embedding similarity: the memory stream of generative agents records timestamped observations and synthesizes higher-order reflections \citep{genagents2023}, while MemGPT introduces an operating-system metaphor in which a fixed context window pages information to and from an external store \citep{memgpt2023}. Subsequent systems impose explicit organization. A-MEM grows a Zettelkasten-style network whose notes are linked and revised as new evidence arrives \citep{amem2025}; AriGraph maintains an integrated semantic--episodic knowledge graph that supports structured planning \citep{arigraph2024}; Zep adds a bi-temporal model that distinguishes event time from ingestion time so that stale facts can be invalidated rather than overwritten \citep{zep2025}; MemoryOS again borrows OS hierarchy to schedule short- and long-term memory \citep{memoryos2025}. Mem0 makes the curation operators first-class, framing maintenance as explicit \textsc{add}/\textsc{update}/\textsc{delete} decisions over extracted facts \citep{mem0_2025}.

Retrieval is the second axis. Dense passage retrieval and the broader retrieval-augmented generation paradigm established the read path that most memory systems still inherit \citep{rag2020}, with late-interaction scoring offering finer token-level matching at scale \citep{colbert2020}. HippoRAG draws on hippocampal indexing to support multi-hop associative recall over a knowledge graph, improving over flat dense retrieval on questions that require chaining facts \citep{hipporag2024}.

Curation quality, not raw capacity, governs whether memory helps. Agentic Context Engineering documents \emph{context collapse}, in which iterative rewriting erodes accumulated detail, and argues for incremental, structure-preserving updates \citep{ace2025}; Mem0's selective edit operators are partly a response to the same failure mode \citep{mem0_2025}; MemoryBank instead imports an Ebbinghaus forgetting curve so that unreinforced traces decay \citep{memorybank2023}. A parallel line exposes memory as an attack surface and a liability: AgentPoison shows that a handful of poisoned entries can hijack downstream retrieval \citep{agentpoison2024}, MemoryGraft demonstrates implanted false memories that persist across sessions \citep{memorygraft2025}, and, even absent an adversary, accumulated memory can \emph{degrade} performance through distraction and error propagation \citep{remembering2026}. This motivates treating memory writes as decisions under reward rather than unconditional logging. Memory-R1 takes exactly this step, using reinforcement learning to assign credit over memory-management actions \citep{memoryr1_2025} --- a \emph{gradient-based} answer to a problem this work attacks training-free.

Evaluation remains immature and almost entirely text-centric. Audits of memory benchmarks reveal that many ostensibly memory-dependent tasks are solvable without genuine cross-session recall \citep{anatomymem2026}, while LongMemEval and LoCoMo provide the most credible long-horizon and conversational probes \citep{longmemeval2024,locomo2024}. On the speech side, retrieval over audio is nascent: WavRAG retrieves directly over acoustic content \citep{wavrag2025}, VoxRAG operates in the spoken domain end-to-end \citep{voxrag2025}, and asynchronous retrieval has been folded into full-duplex spoken dialogue \citep{moshirag2026}. Yet no benchmark isolates \emph{cross-session} memory for a speech agent: the closest, an emotion-in-dialogue suite, is text-derived and single-axis \citep{amber2026}. The gap --- structured, reward-curated, contamination-resistant memory evaluated for spoken interaction --- is precisely where a training-free agent built on a frozen omni model can contribute.

\subsection{Agent Skills and Self-Evolution}
\label{sec:rw-skills}

A complementary form of inference-time accumulation is the \emph{skill}: a reusable procedure the agent acquires and later invokes. A recent systematization formalizes a skill as a tuple of context, policy, transferable trigger, and reward $(C,\pi,T,R)$, separating \emph{when} a skill applies from \emph{what} it does and \emph{how} its utility is scored \citep{sokskills2026}. Voyager pioneered an ever-growing, executable skill library for open-ended control, with new skills written as code and retrieved by description \citep{voyager2023}; agentic workflow memory generalizes this to reusable workflows distilled from successful trajectories \citep{awm2024}; and recent systems automate skill discovery for multimodal use \citep{museautoskill2026}. A central empirical finding tempers the self-generation enthusiasm: on controlled benchmarks, \emph{curated} skill sets substantially outperform self-generated ones, because unverified skills accumulate noise \citep{skillsbench2026}. Two responses follow. Co-evolutionary skill learning trains a surrogate verifier jointly with the skill library so that candidate skills are filtered before admission \citep{coevoskills2026}, and body-first skill routing retrieves by the executable body rather than a possibly-misleading natural-language description \citep{skillrouter2026}. Both can be read as variance-reduction and reward-shaping mechanisms over a discrete skill action space.

Because the target models are frozen, skill optimization here is necessarily gradient-free, aligning with a mature line of prompt-, workflow-, and graph-level optimizers that improve agents without backpropagation: GEPA evolves prompts along a Pareto frontier of reflective mutations \citep{gepa2025}; AFlow searches over workflow graphs \citep{aflow2024}; AgentSquare performs modular agent search \citep{agentsquare2024}; DSPy compiles declarative pipelines \citep{dspy2023}; TextGrad propagates natural-language ``gradients'' \citep{textgrad2024}; ADAS searches the space of agentic designs \citep{adas2024}; and ExpeL extracts and reuses experiential rules \citep{expel2023}. Speech-specific skills are only beginning to appear: rare-word and entity correction via generative error correction \citep{gerrarewords2025}, voice-driven tool use \citep{aura2025}, and paralinguistic style control under explicit reward \citep{paras2s2025}. None of these assembles a frozen omni model's latent competencies into a self-curated, reward-verified skill library --- the construction this work pursues.

\subsection{Agent-Level Convergence and Stability}
\label{sec:rw-convergence}

Iterative self-improvement is not monotone, and a body of evidence catalogues where it stalls or reverses. Reflexion improves through verbal self-feedback but plateaus once reflections cease to surface new information \citep{reflexion2023}. Context engineering collapses when rewriting overwrites detail \citep{ace2025}, and append-only experience logs propagate early errors forward because mistaken entries are never retracted \citep{expfollowing2025}. More starkly, added memory can \emph{hurt}, degrading task accuracy through distraction and compounding of stale facts \citep{remembering2026}. Tree-structured search offers one principled remedy: language-agent tree search couples Monte Carlo rollouts with value backups \citep{lats2023}, inheriting the regret guarantees of UCT bandit exploration \citep{kocsis2006uct}. Convergence under a frozen policy is the explicit object of just-in-time RL, which performs no-gradient credit assignment and characterizes when inference-time self-correction provably converges \citep{jitrl2026}; our analysis treats the same stability question through the lens of a reward-tilted search distribution over a fixed action space, connecting plateau, collapse, and error-propagation phenomena to properties of the tilting reward and its spread $\spread$.

\subsection{Speech Agents and the Open Moat}
\label{sec:rw-speech-agents}

Spoken-language agents are now richly benchmarked, yet almost exclusively in a \emph{single-episode} regime that evaluates a fresh agent on each independent interaction. URO-Bench probes spoken dialogue across understanding, reasoning, and oral conversation \citep{urobench2025}; VoiceBench stress-tests LLM-based voice assistants under realistic perturbations \citep{voicebench2024}; EVA targets expressive and emotional voice interaction \citep{evabench2026}; and a dedicated suite even measures multi-turn memory \emph{within} a single conversation \citep{audiomc2025}. What unifies these is the absence of persistence and self-evolution across sessions: the agent that finishes a dialogue is identical to the one that began it. Where the field does pursue improvement, the mechanism is weight-updating --- continual or alignment training of the speech model itself, as in recent end-to-end conversational systems \citep{cortis2026,soulxduplug2026}. The consequence is an open moat: the \emph{training-free, self-evolving speech agent} --- one that, over its lifetime, accrues reward-curated memory and a verified skill library atop a \emph{frozen} omni model, with no weight or structural change --- is an unfilled slot. This work targets exactly that slot, importing the inference-time reward-tilting and verification machinery of \Cref{sec:rw-memory,sec:rw-skills,sec:rw-convergence} into the spoken-agent setting.

% ===== section 05-theory (writer wordcount 1760) =====
\section{Theory: Optimization-Space Adequacy}
\label{sec:osa}

The thesis of this paper is that training-free reinforcement learning succeeds or fails not because of the cleverness of the search operator, but because of the \emph{adequacy of the action space} over which the search is performed. We make this precise with a sequence of theorems that we have additionally formalized and machine-checked in Lean~4. The development fixes a finite action space $\cZ$ (the inference-time degrees of freedom of a frozen model --- an embedding code in $\cZA$, a generative continuation in $\cZB$, or a product of such), a strictly positive base measure $\qo:\cZ\to\R_{>0}$ with $\sum_{w\in\cZ}\qo(w)=1$ induced by the model at temperature, a bounded reward $\Rwd:\cZ\to\R$, and an inverse-temperature parameter $\beta>0$. All distributions below are probability mass functions on $\cZ$. We write $\Exp_{\qz}[\cdot]$ for expectation under $\qz$ and adopt the convention $0\log 0=0$.

\begin{remark}[Honesty about the formalization]
\label{rem:honesty}
Every result in \Cref{sec:osa} except the two we explicitly flag is \emph{sorry-free} in our Lean development (file \texttt{OptSpace.lean}, supported by the tilting core in \texttt{Tilting.lean}). The two flagged dependencies are isolated as named hypotheses and are \emph{never} discharged silently: (i) Hoeffding's lemma, used only in \Cref{thm:hoeffding}, is taken as an assumption (\texttt{gain\_le\_of\_hoeffding}); and (ii) the order-statistics step in the best-of-$N$ KL bound (the Beirami--Mroueh integral, \texttt{klBoN\_le\_klBoundBoN}) carries the single documented \texttt{sorry} in the wider development and is consumed elsewhere as a hypothesis, not in this section. We state both as assumptions so that no result here over-claims machine verification.
\end{remark}

\subsection{The Tilting Objective and its Optimum}
\label{sec:osa-tilt}

The object of study is the KL-regularized reward functional, whose maximizer is the exponential tilt of the base policy. This is the inference-time specialization of the well-known equivalence between KL-regularized control and Bayesian posterior inference \citep{korbak2022kl}; it underlies controlled decoding \citep{mudgal2024controlled} and the best-of-$N$ literature \citep{beirami2024bon}.

\begin{definition}[Partition function, Gibbs optimum, regularized objective]
\label{def:tilt}
Define the \emph{partition function}
\[
\Zpart \;:=\; \sum_{w\in\cZ} \qo(w)\,\exp\!\big(\Rwd(w)/\beta\big)\;=\;\Exp_{\qo}\!\big[\exp(\Rwd/\beta)\big]\;>\;0,
\]
the \emph{Gibbs optimum} (exponential tilt)
\[
\qs(z) \;:=\; \frac{\qo(z)\,\exp\!\big(\Rwd(z)/\beta\big)}{\Zpart},
\]
which is a probability distribution because $\qo>0$ and $\Zpart>0$, and the \emph{KL-regularized objective}
\[
\Fobj{\qz} \;:=\; \sum_{z\in\cZ}\qz(z)\,\Rwd(z) \;-\; \beta\sum_{z\in\cZ}\qz(z)\,\log\frac{\qz(z)}{\qo(z)} \;=\; \Exp_{\qz}[\Rwd] \;-\; \beta\,\KL{\qz}{\qo}.
\]
\end{definition}

The next lemma is the algebraic spine of the whole section: it converts the gap to the optimum into a single nonnegative KL divergence.

\begin{lemma}[Master identity]
\label{lem:master}
For every probability distribution $\qz$ on $\cZ$,
\[
\Fobj{\qz} \;=\; \beta\log\Zpart \;-\; \beta\,\KL{\qz}{\qs},
\qquad\text{hence}\qquad
\Fobj{\qs}-\Fobj{\qz} \;=\; \beta\,\KL{\qz}{\qs}.
\]
\end{lemma}

\begin{proof}
Taking logarithms in \Cref{def:tilt}, $\log\qs(z)=\log\qo(z)+\Rwd(z)/\beta-\log\Zpart$, so that, pointwise,
\[
\Rwd(z) \;=\; \beta\Big(\log\tfrac{\qs(z)}{\qo(z)} + \log\Zpart\Big).
\]
Substituting into $\Fobj{\qz}$ and using $\sum_z\qz(z)=1$,
\begin{align*}
\Fobj{\qz}
&= \sum_z \qz(z)\Big[\beta\log\tfrac{\qs(z)}{\qo(z)} + \beta\log\Zpart\Big] - \beta\sum_z \qz(z)\log\tfrac{\qz(z)}{\qo(z)} \\
&= \beta\log\Zpart + \beta\sum_z \qz(z)\Big[\log\tfrac{\qs(z)}{\qo(z)} - \log\tfrac{\qz(z)}{\qo(z)}\Big]
 = \beta\log\Zpart + \beta\sum_z \qz(z)\log\tfrac{\qs(z)}{\qz(z)} \\
&= \beta\log\Zpart - \beta\,\KL{\qz}{\qs}.
\end{align*}
Evaluating at $\qz=\qs$ gives $\Fobj{\qs}=\beta\log\Zpart$ (the KL term vanishes), and subtracting yields the stated gap.
\end{proof}

\begin{theorem}[Tilting optimality; \texttt{tilting\_optimal}, T1]
\label{thm:t1}
For every probability distribution $\qz$ on $\cZ$, $\Fobj{\qz}\le \Fobj{\qs}$, with equality iff $\qz=\qs$. Thus the exponential tilt $\qs$ is the unique maximizer of the inference-time objective $\Fobj{\cdot}$.
\end{theorem}

\begin{proof}
By \Cref{lem:master}, $\Fobj{\qs}-\Fobj{\qz}=\beta\,\KL{\qz}{\qs}$. Nonnegativity of KL (Gibbs' inequality) follows from $\log x\le x-1$ for $x>0$: writing the ratio $r(z)=\qs(z)/\qz(z)$ on the support of $\qz$,
\[
-\KL{\qz}{\qs}=\sum_z \qz(z)\log r(z)\;\le\;\sum_z \qz(z)\big(r(z)-1\big)=\sum_z \qs(z)-\sum_z \qz(z)=0,
\]
so $\KL{\qz}{\qs}\ge0$ and $\Fobj{\qz}\le\Fobj{\qs}$. Since $\beta>0$ and $\log x=x-1$ only at $x=1$, equality forces $\qz=\qs$ pointwise.
\end{proof}

\Cref{thm:t1} is the entire ceiling of any training-free method: no reranker, decoder, or search distribution $\qz\in\Delta(\cZ)$ can exceed $\Fobj{\qs}=\beta\log\Zpart$. The rest of the section asks the operational question this raises --- \emph{how large is the headroom $\Fobj{\qs}-\Fobj{\qo}$ between the model's own sampling distribution $\qo$ and this ceiling?} --- because that headroom is exactly the gain a training-free optimizer can ever realize.

\subsection{OSA-1: A Flat or Degenerate Action Space Yields Vanishing Gain}
\label{sec:osa1}

\begin{definition}[Optimization gain]
\label{def:gain}
The \emph{gain} of the action space under reward $\Rwd$ at temperature $\beta$ is the realizable headroom $\gain := \Fobj{\qs}-\Fobj{\qo}$.
\end{definition}

\begin{proposition}[Closed form for the gain; \texttt{gain\_eq}, OSA-1a]
\label{prop:gaineq}
\[
\gain \;=\; \beta\log\Zpart - \Exp_{\qo}[\Rwd] \;=\; \beta\log\Exp_{\qo}\!\big[\exp(\Rwd/\beta)\big] - \Exp_{\qo}[\Rwd].
\]
\end{proposition}

\begin{proof}
Apply \Cref{lem:master} at $\qz=\qo$. Since $\KL{\qo}{\qo}=0$, the definition gives $\Fobj{\qo}=\Exp_{\qo}[\Rwd]-\beta\KL{\qo}{\qo}=\Exp_{\qo}[\Rwd]$, while $\Fobj{\qs}=\beta\log\Zpart$. Subtract, and recall $\Zpart=\Exp_{\qo}[\exp(\Rwd/\beta)]$ from \Cref{def:tilt}.
\end{proof}

\begin{proposition}[Nonnegativity of the gain; \texttt{gain\_nonneg}, OSA-1b]
\label{prop:gainnn}
$\gain\ge0$, with equality iff $\qs=\qo$.
\end{proposition}

\begin{proof}
Two equivalent routes. (i) By \Cref{lem:master}, $\gain=\Fobj{\qs}-\Fobj{\qo}=\beta\,\KL{\qo}{\qs}\ge0$, strict unless $\qo=\qs$. (ii) Directly from \Cref{prop:gaineq}, Jensen's inequality applied to the strictly convex $t\mapsto e^{t/\beta}$ gives $\beta\log\Exp_{\qo}[e^{\Rwd/\beta}]\ge \beta\,\Exp_{\qo}[\log e^{\Rwd/\beta}]=\Exp_{\qo}[\Rwd]$.
\end{proof}

The gain is therefore a divergence between what the model already does, $\qo$, and what the reward would have it do, $\qs$. When the action space is \emph{degenerate} --- so that the reward cannot distinguish its elements --- this divergence collapses.

\begin{theorem}[Flat reward, no gain; \texttt{flat\_no\_gain}, OSA-1c $\Rightarrow$ T3]
\label{thm:flat}
If $\Rwd\equiv c$ is constant on $\cZ$, then $\qs=\qo$ and $\gain=0$.
\end{theorem}

\begin{proof}
With $\Rwd(z)=c$, $\Zpart=\sum_w\qo(w)e^{c/\beta}=e^{c/\beta}$, so $\qs(z)=\qo(z)e^{c/\beta}/e^{c/\beta}=\qo(z)$ (this is exactly T3, \texttt{qstar\_eq\_q0\_of\_const}). By \Cref{prop:gaineq}, $\gain=\beta\log e^{c/\beta}-c=c-c=0$.
\end{proof}

\Cref{thm:flat} is the formal content of our central negative observation about \emph{single-model, output-as-context} training-free RL. When the only inference-time degree of freedom is the model's own decoded continuation fed back as its next context --- with a reward that is near-constant across the plausible continuations the model actually samples, as is empirically the case for self-consistency-style signals on a frozen model with no external verifier --- the effective $\Rwd$ is approximately flat on $\mathrm{supp}(\qo)$, the tilt does essentially nothing, and there is no leverage for the optimizer to exploit. This is consistent with reports that purely self-referential test-time RL plateaus quickly \citep{zuo2025ttrl,reflexion2023} and with spurious-reward phenomenology \citep{spuriousrewards2025,skalse2022rewardhacking}. The next theorem makes the dependence on reward dispersion quantitative.

\begin{assumption}[Hoeffding's lemma; isolated, cf.\ \Cref{rem:honesty}]
\label{ass:hoeffding}
Let $X$ be a random variable supported in an interval of length $\ell$ a.s.\ under a measure $\mu$. Then for all $\lambda\in\R$,
\[
\log\Exp_{\mu}\!\big[e^{\lambda X}\big]\;\le\;\lambda\,\Exp_{\mu}[X]+\frac{\lambda^2\ell^2}{8}.
\]
In Lean this is the hypothesis named in \texttt{gain\_le\_of\_hoeffding}; we do not prove it from first principles in our development.
\end{assumption}

\begin{theorem}[Quantitative collapse of gain; \texttt{gain\_le\_of\_hoeffding}, OSA-1d]
\label{thm:hoeffding}
Let $\spread:=\max_{z\in\cZ}\Rwd(z)-\min_{z\in\cZ}\Rwd(z)$ be the reward spread on $\cZ$. Under \Cref{ass:hoeffding},
\[
\gain \;\le\; \frac{\spread^{2}}{8\beta}.
\]
In particular $\spread\to0\Rightarrow\gain\to0$: a vanishing reward spread forces a vanishing realizable gain, uniformly in the search operator.
\end{theorem}

\begin{proof}
Apply \Cref{ass:hoeffding} with $\mu=\qo$, $X=\Rwd$ (supported in an interval of length $\spread$ on $\cZ$), and $\lambda=1/\beta$:
\[
\log\Exp_{\qo}\!\big[e^{\Rwd/\beta}\big]\;\le\;\frac1\beta\,\Exp_{\qo}[\Rwd]+\frac{1}{\beta^2}\cdot\frac{\spread^2}{8}.
\]
Multiply by $\beta>0$ and rearrange to $\beta\log\Exp_{\qo}[e^{\Rwd/\beta}]-\Exp_{\qo}[\Rwd]\le \spread^2/(8\beta)$, which is exactly $\gain$ by \Cref{prop:gaineq}. The limit is immediate.
\end{proof}

\Cref{thm:hoeffding} pairs with \Cref{prop:gainnn} to sandwich the gain, $0\le\gain\le\spread^2/(8\beta)$, and identifies the reward spread on the realized support as the sole lever of a frozen model: it is the inference-time analogue of the reward-model scaling laws that bound over-optimization \citep{gao2023overopt}, and it explains why best-of-$N$ and reranking deliver returns proportional to how much the verifiable reward actually separates candidates \citep{beirami2024bon,khalaf2025hedgetune,snell2024ttscompute}.

\subsection{OSA-2: Context Isolation Enlarges the Space --- Additive, Growing Gain}
\label{sec:osa2}

\Cref{thm:flat,thm:hoeffding} are a diagnosis, not a verdict: they say that a \emph{degenerate} space yields no gain, leaving open whether the space can be enlarged. It can. The mechanism is \emph{context isolation}: factoring a monolithic decision into components that are decoded under separate, non-interfering contexts (distinct sub-agents, or skill/memory-conditioned roles), each carrying its own verifiable reward. Formally this replaces $\cZ$ by a product and $\Rwd$ by a separable sum.

\begin{lemma}[Partition function factorizes; \texttt{Zpart\_product}]
\label{lem:zfactor}
Let $\cZ=\cZ_1\times\cZ_2$ with base $\qo=q_0^{(1)}\otimes q_0^{(2)}$ and \emph{separable} reward $\Rwd(z_1,z_2)=\Rwd_1(z_1)+\Rwd_2(z_2)$ (context isolation). Then $\Zpart=\Zpart^{(1)}\,\Zpart^{(2)}$, where $\Zpart^{(i)}=\sum_{z_i}q_0^{(i)}(z_i)e^{\Rwd_i(z_i)/\beta}$.
\end{lemma}

\begin{proof}
A direct computation using independence and the multiplicativity of the exponential over the separable reward:
\[
\Zpart=\sum_{z_1,z_2}q_0^{(1)}(z_1)q_0^{(2)}(z_2)\,e^{(\Rwd_1(z_1)+\Rwd_2(z_2))/\beta}
=\Big(\sum_{z_1}q_0^{(1)}(z_1)e^{\Rwd_1(z_1)/\beta}\Big)\Big(\sum_{z_2}q_0^{(2)}(z_2)e^{\Rwd_2(z_2)/\beta}\Big).\qedhere
\]
\end{proof}

\begin{theorem}[Gain is additive over isolated components; \texttt{gain\_product}, OSA-2]
\label{thm:gainprod}
Under the hypotheses of \Cref{lem:zfactor}, $\gain=\gain(\Rwd_1)+\gain(\Rwd_2)$, where $\gain(\Rwd_i)=\beta\log\Zpart^{(i)}-\Exp_{q_0^{(i)}}[\Rwd_i]$. Consequently, for $k$ context-isolated components with product base and additively separable reward $\Rwd=\sum_{i=1}^k\Rwd_i$,
\[
\gain=\sum_{i=1}^{k}\gain(\Rwd_i),
\]
and this sum is \emph{strictly increasing in $k$} whenever each newly added component $\Rwd_i$ is non-constant on its support, since by \Cref{prop:gainnn} each summand satisfies $\gain(\Rwd_i)>0$ exactly when $\Rwd_i$ is non-constant.
\end{theorem}

\begin{proof}
By \Cref{prop:gaineq} and \Cref{lem:zfactor}, $\beta\log\Zpart=\beta\log\Zpart^{(1)}+\beta\log\Zpart^{(2)}$. Under the product base with separable reward, $\Exp_{\qo}[\Rwd]=\Exp_{q_0^{(1)}}[\Rwd_1]+\Exp_{q_0^{(2)}}[\Rwd_2]$. Subtracting term by term gives the two-factor claim; the $k$-fold statement follows by induction on the components. Strict positivity of a summand is the equality case of \Cref{prop:gainnn} (equivalently, strict Jensen in \Cref{prop:gaineq}): $\gain(\Rwd_i)=0$ iff $\Rwd_i$ is constant.
\end{proof}

\Cref{thm:gainprod} is the constructive counterpart of \Cref{sec:osa1}. Where a single near-flat space affords no leverage, \emph{enlarging} the action space by context isolation --- decomposing the task across sub-agents and loading distinct skills/memories so that each carries an independent, non-degenerate verifiable reward --- turns a degenerate problem into a product problem whose gain is the sum of per-component gains and grows with the number of isolated, reward-bearing components. This is the mechanism by which agentic decomposition, retrieval-conditioned roles, and verifiable per-skill rewards \citep{tulu3,jitrl2026} convert a setting where training-free RL is impotent into one where it is attainable.

\subsection{OSA-3: Rollout Instability and Credit-Assigned Convergence}
\label{sec:osa3}

Enlarging the space (OSA-2) makes a positive gain \emph{available}; OSA-3 concerns whether a rollout policy can \emph{reach} it. The master identity already settles the geometry.

\begin{theorem}[Rollout deficit; \texttt{rollout\_deficit}, OSA-3a]
\label{thm:rollout}
Every rollout policy $\qz$ on $\cZ$ satisfies
\[
\Fobj{\qz} \;=\; \Fobj{\qs} \;-\; \beta\,\KL{\qz}{\qs},\qquad \beta\,\KL{\qz}{\qs}\ge0,
\]
with the deficit $\beta\,\KL{\qz}{\qs}$ strictly positive unless $\qz=\qs$. Hence a naive rollout that fails to land on $\qs$ incurs a strictly positive, irreducible objective deficit, and there is no \emph{a priori} monotone path in $\qz$-space that decreases this deficit to zero.
\end{theorem}

\begin{proof}
This is \Cref{lem:master} read as a statement about an arbitrary rollout $\qz$; strictness is the equality case of \Cref{thm:t1}. The KL surface $\qz\mapsto\KL{\qz}{\qs}$ is convex with its unique zero at $\qs$, so a sequence of rollout updates not steered toward $\qs$ need not be monotone in the deficit.
\end{proof}

\begin{theorem}[Per-component tilting reaches the global optimum; \texttt{qstar\_product}, OSA-3b]
\label{thm:qstarprod}
Under the isolated product structure of \Cref{lem:zfactor}, the global Gibbs optimum factorizes:
\[
\qs(z_1,z_2)\;=\;q^{*}_{1}(z_1)\,q^{*}_{2}(z_2),\qquad q^{*}_{i}(z_i)=\frac{q_0^{(i)}(z_i)\,e^{\Rwd_i(z_i)/\beta}}{\Zpart^{(i)}}.
\]
Therefore tilting each isolated, credit-assigned component to its own local optimum $q^{*}_i$ assembles \emph{exactly} the global optimum $\qs$, which by \Cref{thm:t1} is $\Fobj{\cdot}$-optimal over all of $\Delta(\cZ)$.
\end{theorem}

\begin{proof}
Using $\qo=q_0^{(1)}\otimes q_0^{(2)}$, the separable reward, and \Cref{lem:zfactor},
\[
\qs(z_1,z_2)=\frac{q_0^{(1)}(z_1)q_0^{(2)}(z_2)\,e^{\Rwd_1(z_1)/\beta}e^{\Rwd_2(z_2)/\beta}}{\Zpart^{(1)}\Zpart^{(2)}}
=\frac{q_0^{(1)}(z_1)e^{\Rwd_1(z_1)/\beta}}{\Zpart^{(1)}}\cdot\frac{q_0^{(2)}(z_2)e^{\Rwd_2(z_2)/\beta}}{\Zpart^{(2)}}=q^{*}_1(z_1)\,q^{*}_2(z_2).
\]
Global optimality is \Cref{thm:t1}.
\end{proof}

\paragraph{The trust-region hinge.}
\Cref{thm:rollout,thm:qstarprod} together specify when training-free RL converges. The $\beta\,\KL{\qz}{\qs}$ term in \Cref{lem:master} is not cosmetic regularization: it is a \emph{trust region} that penalizes drift from the tilted optimum and thereby enforces the slow-drift precondition under which a retrieval-/advantage-tilting update --- such as JitRL's no-gradient credit-assigned tilt \citep{jitrl2026} --- contracts toward $\qs$. Concretely, \Cref{thm:qstarprod} shows that performing the tilt \emph{per isolated component} (credit assignment) is not a heuristic shortcut but reaches the same global optimum as a monolithic tilt, while \Cref{thm:rollout} shows that any policy off $\qs$ retains a strictly positive deficit. The hinge linking the two is the temperature $\beta$: a large enough trust region keeps each rollout in the basin where per-component tilting is a descent on $\KL{\cdot}{\qs}$ (OSA-3a's non-convergence becomes OSA-3b's convergence), whereas a fast-drift rollout that ignores the trust region --- appending its own unverified outputs as context and chasing them --- has no monotone guarantee and is observed to diverge through context collapse and error propagation \citep{reflexion2023,ace2025,gao2023overopt}. In sum, OSA-1 says a degenerate space cannot be optimized, OSA-2 says context isolation enlarges it into one that can, and OSA-3 says the resulting gain is reachable precisely when credit-assigned tilting is run inside the KL trust region.

% ===== section 06-convergence (writer wordcount 1180) =====
\section{Convergence Analysis}
\label{sec:convergence}

A central claim of this work is that the inference-time degrees of freedom $\cZ$ admit a sharp asymmetry: convergence is \emph{provably} controllable at the output level, where the search distribution acts directly on a fixed generative operator, yet only \emph{asymptotically} controllable at the agent level, where memory, skills, and context co-evolve with the policy. This section surveys and synthesizes what is established, separates proven finite-$N$ guarantees from empirical or absent ones, and isolates the open problem that the Operator-Search Architecture (OSA) frames. Throughout we measure convergence against the Gibbs optimum $\qs \propto \qo \exp(\Rwd/\beta)$ of the KL-regularized objective $\Fobj{q} = \Exp_{q}[\Rwd] - \beta\,\KL{q}{\qo}$, whose maximizer is exactly $\qs$ with optimal value $\beta \log \Zpart$ \citep{korbak2022kl}. The question ``does method $M$ converge'' is then made precise as: does the realized $q$ approach $\qs$ (or attain $\Fobj{\qs}$) as the compute budget $N$ grows, and at what rate?

\subsection{Output-Level: Proven Finite-N Convergence}
\label{subsec:output-level}

At the output level the action lives in $\cZB$ (and, for reranking, in $\cZA$): the generative operator is frozen and the only freedom is which candidates to draw, weight, or select. Here the literature provides genuine non-asymptotic guarantees.

\paragraph{Soft best-of-$N$ tilts to $\qs$ at rate $O(1/N)$.} The cleanest result is for soft (regularized) best-of-$N$: drawing $N$ i.i.d.\ samples from $\qo$ and resampling one with weights $\propto \exp(\Rwd/\beta)$ yields a distribution whose KL to the Gibbs optimum $\qs$ decays as $O(1/N)$, so the induced $q$ is a consistent, rate-quantified estimator of $\qs$ \citep{verdun2025softbon}. This is the tightest convergence statement available for any element of $\cZ$ and serves as the reference behavior the rest of the architecture is measured against.

\paragraph{MBR and selection-style decoders.} Minimum Bayes risk decoding converges to the risk-minimizing output at the Monte-Carlo rate $O(n^{-1/2})$ in the number of pseudo-references \citep{ichihara2025mbr}, a guarantee that transfers directly to ASR, where MBR over hypothesis lists reduces word error in a way that is consistent rather than merely heuristic \citep{mbrasr2025}. Hybrids that fuse MBR weighting with best-of-$N$ selection inherit both the variance reduction of the former and the reward-alignment of the latter \citep{jinnai2024mbrbon}, while self-consistency is the special case of majority-vote MBR under a $0/1$ kernel \citep{wang2023selfconsistency}.

\paragraph{Hard best-of-$N$: KL budget and regret.} Hard best-of-$N$ does not target $\qs$ but is nonetheless tightly characterized. Its KL to the base distribution is bounded by $\log N - (N-1)/N$ \citep{beirami2024bon}, the inequality we formalize as Lean theorem T2; this caps how far selection can move the output distribution and hence bounds reward inflation. Complementarily, a smoothing-lens analysis shows the suboptimality (regret) of best-of-$N$ against the optimal tilt grows only as $O(\sqrt{\log N})$ \citep{aminian2025smoothing}, our Lean theorem T6. Together these bracket the operator: selection buys reward at a logarithmically growing KL price and a sub-polynomially growing regret, both finite and explicit. Related asymptotic-equivalence results connect best-of-$N$ to KL-regularized RL in the large-$N$ limit \citep{yang2024asymptotics,gui2024bonbon}.

\paragraph{The over-optimization hump and the $N^\ast$ controller.} Convergence \emph{toward $\qs$} is not the same as convergence \emph{of measured reward toward true utility}. When $\Rwd$ is a proxy, pushing $N$ upward eventually over-optimizes: true performance rises, peaks, and declines, tracing a hump whose location scales predictably with the KL budget \citep{gao2023overopt}. This motivates treating $N$ (equivalently the effective temperature) as a controlled variable. HedgeTune computes the optimal $N^\ast$ — the best-of-Poisson count or tuning point that maximizes realized utility before the proxy diverges from truth — turning the budget into a closed-form set-point rather than a hyperparameter \citep{khalaf2025hedgetune}. The existence of a finite $N^\ast$ is itself a convergence statement: more compute past $N^\ast$ is provably counterproductive under a fixed proxy.

\paragraph{Compute allocation and verifier-guided variants.} Several methods sharpen the constant rather than the rate. Guided speculative inference composes a small proposal with a large verifier to reach the tilted target at lower cost \citep{geuter2025gsi}; CarBoN calibrates per-instance temperature to allocate samples where the reward landscape is informative \citep{carbon2025}; certified self-consistency and confidence-guided escalation give stopping rules with coverage guarantees \citep{certifiedsc2025,cges2025}; and test-time-compute-optimal scaling shows the budget allocation itself can be optimized \citep{snell2024ttscompute}. Sequential Monte-Carlo and Feynman--Kac steering extend the same tilting calculus to token-level resampling with consistency under standard SMC assumptions \citep{zhao2024twistedsmc,singhal2025fksteering}. The common thread: every result in this subsection concerns a \emph{fixed} operator acted on by a search distribution, which is precisely why finite-$N$ analysis is tractable.

\subsection{Agent-Level: Asymptotic Consistency Under a Trust Region}
\label{subsec:agent-level}

Moving from $\cZB$ to the full memory$\cdot$skill$\cdot$context space changes the operator itself between steps, and the guarantees weaken from finite-$N$ to asymptotic.

\paragraph{JitRL: an exact KL-constrained update, consistent in probability.} The strongest agent-level result is JitRL's no-gradient credit assignment \citep{jitrl2026}. It proves that the embedding update $z' = z + \beta \hat{A}$ realizes \emph{exactly} the KL-constrained optimal policy $\pi^\ast \propto \pi_\theta \exp(\beta A)$ — i.e.\ the same Gibbs tilt as \Cref{subsec:output-level}, now applied in $\cZA$ to a frozen backbone. The advantage $A$ is unknown and is estimated by a $k$-nearest-neighbor estimator $\hat{A}$; this estimator is consistent \emph{in probability} only under a slow-policy-drift precondition: the neighborhood statistics must be drawn under a policy that changes slowly enough for the local estimate to remain valid. The resulting guarantee is therefore \emph{asymptotic} (the realized policy approaches $\pi^\ast$ as the sample pool grows under bounded drift), \emph{not} a finite-time monotone-improvement statement. The trust region is implicit in the drift condition: violate it and the consistency argument lapses with no replacement bound. This is the central honest caveat — JitRL converges to the right object, but only in the limit and only inside a region the analysis assumes rather than enforces.

\paragraph{Tree search lacks transferable optimality.} Agent planners that wrap a language model in Monte-Carlo tree search, such as LATS \citep{lats2023}, borrow the UCT selection rule whose regret guarantees were proven for stationary, bounded-reward bandits \citep{kocsis2006uct}. Those guarantees do not transfer: the per-node reward here is a non-stationary, self-generated value estimate, the action set is unbounded, and the ``environment'' is the model's own context, violating the i.i.d.\ and stationarity premises UCT requires. Empirically LATS improves, but its convergence is not certified by the theory it invokes. Variance-reduced policy-gradient estimators such as RLOO \citep{ahmadian2024rloo,greensmith2004variance} face the analogous gap once the operator is non-stationary.

\subsection{The Gap and the Failure Modes}
\label{subsec:gap}

The asymmetry is not merely a proof-technique artifact: naive rollout in the agent space is observed to \emph{not} converge along every action axis, which is what makes the missing guarantee consequential.

\paragraph{Documented non-convergence across all axes.} On the context/reflection axis, iterated self-refinement plateaus — Reflexion-style loops stop improving and oscillate after a few rounds \citep{reflexion2023}. On the context-construction axis, accumulated context collapses, degrading rather than enriching the operator \citep{ace2025}. On the memory axis, append-only experience logs propagate early errors forward, and added memory can actively hurt downstream accuracy \citep{expfollowing2025,remembering2026}. These are not isolated anecdotes but a coherent pattern: when the search distribution is allowed to rewrite the operator without a trust region, the very Gibbs-tilt logic that converges at the output level becomes a divergent feedback loop, because the reference $\qo$ drifts under the update. Proxy-reward exploitation compounds this \citep{skalse2022rewardhacking,gao2023overopt}.

\paragraph{What is and is not proven.} \Cref{tab:convergence} summarizes the landscape. Proven finite-$N$ convergence exists \emph{only} at the output level, where the operator is fixed. At the agent level the best available result is asymptotic consistency conditioned on slow drift; everything else is empirical or has no transferable guarantee. The open problem OSA frames is therefore precise: a \emph{finite-time, monotone-improvement} guarantee over the full $\cZ = \cZA \cup \cZB$ space — a single trust region, analogous to the drift condition but \emph{enforced} rather than assumed, under which each update to memory, skills, context, or decoding is certified not to decrease realized utility. The output-level results give the template (a Gibbs tilt with a KL budget and a tuned $N^\ast$); the agent-level results give the obstruction (a moving reference and a non-stationary advertised reward). Closing the gap means transporting the finite-$N$ guarantees of \Cref{subsec:output-level} across an operator that is itself part of the action.

\begin{table}[t]
\centering
\small
\begin{tabular}{@{}lll@{}}
\toprule
Method & Convergence & Scope \\
\midrule
Soft best-of-$N$ \citep{verdun2025softbon} & Proven, $O(1/N)$ to $\qs$ & Output ($\cZB$) \\
MBR decoding \citep{ichihara2025mbr,mbrasr2025} & Proven, $O(n^{-1/2})$ & Output ($\cZB$) \\
Best-of-$N$ KL bound \citep{beirami2024bon} & Proven (Lean T2) & Output ($\cZB$) \\
Best-of-$N$ regret \citep{aminian2025smoothing} & Proven, $O(\sqrt{\log N})$ (Lean T6) & Output ($\cZB$) \\
$N^\ast$ controller \citep{khalaf2025hedgetune} & Proven set-point & Output ($\cZB$) \\
Over-optimization hump \citep{gao2023overopt} & Proven scaling law & Output ($\cZB$) \\
GSI / CarBoN \citep{geuter2025gsi,carbon2025} & Proven (target-consistent) & Output ($\cZB$) \\
Self-consistency \citep{wang2023selfconsistency} & Proven (MBR special case) & Output ($\cZB$) \\
\midrule
JitRL \citep{jitrl2026} & Asymptotic, slow-drift precond. & Embedding ($\cZA$) \\
LATS / UCT \citep{lats2023,kocsis2006uct} & None transferable & Agent (plan) \\
\midrule
Reflexion loop \citep{reflexion2023} & None (plateau) & Context \\
ACE context \citep{ace2025} & None (collapse) & Context \\
Append-only memory \citep{expfollowing2025} & None (error propagation) & Memory \\
Memory accumulation \citep{remembering2026} & None (contamination) & Memory \\
\bottomrule
\end{tabular}
\caption{Convergence status of inference-time methods across the action space $\cZ$. Finite-$N$ guarantees exist only where the generative operator is frozen; agent-level methods are at best asymptotic under an assumed trust region, and naive rollout over memory/context is documented to diverge.}
\label{tab:convergence}
\end{table}

% ===== section 07-feasibility (writer wordcount 1268) =====
\section{Feasibility Analysis}
\label{sec:feasibility}

We now argue that the agent-level program of this work is not only well-posed but already \emph{instantiable} from off-the-shelf components, and that our own preliminary measurements both motivate and constrain the design. The argument proceeds in four steps: we identify the two classes of frozen omni models with the two structural roles an episodic agent requires (\Cref{sec:feas:components}); we show that a single verifiable-reward acceptance gate furnishes the control law shared by memory and skills (\Cref{sec:feas:gate}); we state an honest novelty delta against a crowded mechanism literature (\Cref{sec:feas:novelty}); and we present in-house results that locate the operating regime empirically (\Cref{sec:feas:prelim}).

\subsection{The Two Omni Classes as Agent Components}
\label{sec:feas:components}

The decisive enabling observation is that the contemporary omni-model ecosystem ships, off the shelf, the two computational primitives an episodic agent needs, and that both can be used \emph{frozen}. On one side, \emph{vector-omni bi-encoders} map an audio (or audio--text) input to a fixed embedding in a shared retrieval space \citep{omniembed2025,omniembedaudio2026}; their natural role is the agent's \emph{memory and index}, the read/write substrate over past episodes, exactly as speech-native retrieval-augmented systems already use them \citep{wavrag2025,speechrag2024,voxrag2025}. On the other side, \emph{generative thinker--talker} omni models map an input to a token sequence under a decoding policy \citep{qwen3omni,qwen25omni,salmonn2023}; their natural role is the agent's \emph{policy and skill executor}. Both classes are pretrained, publicly released, and, for our purposes, \emph{frozen}: the action lives entirely in the inference-time degrees of freedom $\cZ$, partitioned into the embedding/retrieval operator $\cZA$ realised by the bi-encoder and the generative operator $\cZB$ realised by the thinker--talker. This separation is not a convenience but a structural commitment: it lets us treat memory construction (Operator-A) and skill execution (Operator-B) as two faces of a single search distribution $\qz$ over $\cZ$, optimised against the same reward $\Rwd$ without touching weights or architecture. Crucially, both primitives already exist at the quality frontier, so feasibility reduces to whether a \emph{control law} can compose them productively rather than to whether the parts can be built.

\subsection{The Verifiable-Reward Acceptance Gate as the Control Law}
\label{sec:feas:gate}

The control law is a verifiable-reward \emph{acceptance gate}: a candidate memory write or skill invocation is admitted only if it improves a verifiable reward on held-out probes, and is otherwise rejected. This single mechanism governs both faces of $\cZ$. The case for making the gate the load-bearing component is empirical and, in our reading, decisive. \textsc{SkillsBench} reports that \emph{curated} skills raise agent success by $+16.2$ percentage points, whereas \emph{self-generated} skills contribute essentially zero in aggregate, with $16$ of $84$ skills measurably \emph{negative} \citep{skillsbench2026}. The lesson is not that self-evolution is futile but that unfiltered self-evolution is: without an acceptance test, an append-only memory or skill library accrues as much harm as benefit, a failure mode echoed across the agent-stability literature \citep{expfollowing2025,ace2025,remembering2026,reflexion2023}. The gate is precisely the missing filter. When ground-truth labels are scarce, a surrogate verifier can stand in for the reward without collapsing the guarantee, as co-evolved skill--verifier pairs demonstrate \citep{coevoskills2026}; and for speech specifically the verifiable rewards already exist, from paralinguistic-aware speech reward models \citep{gsrm2026,paras2s2025} to the WER/exact-match/intent rewards reused throughout this series. We therefore adopt the gate as the one control law shared by $\cZA$ (admit a memory entry) and $\cZB$ (admit a skill or decoding action), with the surrogate verifier as the fallback when labelled probes run out. This is the same KL-regularised tilting that underwrites best-of-$N$ and reward-guided decoding \citep{korbak2022kl,mudgal2024controlled,beirami2024bon}, here lifted from token decisions to episode-level admission decisions.

\subsection{Novelty Delta}
\label{sec:feas:novelty}

We state the contribution honestly. The \emph{mechanism}---self-evolving episodic memory plus a growing skill library refined by feedback---is crowded: reflective self-improvement \citep{reflexion2023,expel2023}, library-building agents \citep{voyager2023,awm2024,museautoskill2026}, and gradient-free inference-time credit assignment \citep{jitrl2026} all occupy adjacent ground. Our claim is not the mechanism in the abstract. The unfilled slot is the specific composition: a \emph{frozen omni policy} (Operator-B) and a \emph{frozen omni memory} (Operator-A), driven by \emph{episodic self-evolution} under a \emph{verifiable speech reward}, with no weight updates anywhere in the loop. The closest priors fall short on exactly this axis. Trainable instruction embedders such as \textsc{bge-en-icl} require gradient training to acquire their in-context behaviour \citep{bgeenicl2024}, violating the frozen-memory constraint; few-shot--emergent audio models such as \textsc{MiMo-Audio} obtain their capability through pretraining rather than through inference-time activation of an already-frozen model \citep{mimoaudio2025}. Neither realises the all-frozen, reward-gated, speech-verifiable agent loop. To our knowledge, this composition has not been demonstrated, and we make the corresponding ``first, to our knowledge'' claim while being explicit that each ingredient in isolation has prior art.

\subsection{Preliminary Empirical Evidence}
\label{sec:feas:prelim}

Our in-house measurements locate the operating regime and, in doing so, supply the empirical motivation for moving the action up from single-model conditioning to the agent action space. All numbers below are our own preliminary results on a frozen \texttt{omni-embed-nemotron-3b}, with kNN probes, seed $42$, and $1000$-sample bootstrap confidence intervals (CREMA-D dev$600$/test$300$; MInDS-14 dev$280$/test$257$); they are reported here verbatim and labelled as preliminary.

\paragraph{Single-model conditioning is near-flat (OSA-1 in practice).}
\Cref{tab:cremad-matrix} reports test probe accuracy for the frozen embedding under four instruction-conditioning regimes, crossed against three readout factors. The same frozen embedding simultaneously supports near-perfect content readout ($\approx 1.00$), weak emotion readout ($\approx 0.36$), and at-chance speaker readout ($\approx 0.04$). Decisively, \emph{instruction conditioning does not move the factor accuracies}: the emotion delta from baseline to emotion-instruction is $+0.007$ with overlapping CIs, the speaker delta is $-0.003$, and the conditioning matrix is not diagonally dominant. The picture is stable across the architecture: a sweep over all $37$ hidden states and two pooling readouts leaves speaker accuracy at chance at \emph{every} layer (full-sequence best $0.043$ at L36; audio-token best $0.043$ at L32), with content saturating early (full-seq $0.990$ at L24) and emotion peaking only weakly (full-seq $0.400$ at L0). This flatness is exactly the empirical signature predicted by the near-flat single-model action space (the OSA-1 condition): within one frozen model, instruction-conditioning offers essentially no reward spread $\spread$ to optimise over, so the gain $\gain$ from any Gibbs-tilted search $\qz$ over single-model conditioning is bounded near zero.

\begin{table}[t]
\centering
\caption{In-house preliminary results. CREMA-D conditioning$\times$factor TEST probe accuracy (frozen \texttt{omni-embed-nemotron-3b}, seed $42$). Instruction conditioning is flat across all factors (\texttt{diagonal\_dominant}$=$False).}
\label{tab:cremad-matrix}
\begin{tabular}{lccc}
\toprule
Conditioning & Content & Emotion & Speaker \\
\midrule
Baseline              & 0.997 & 0.360 & 0.040 \\
Content-instruction   & 1.000 & 0.333 & 0.043 \\
Emotion-instruction   & 1.000 & 0.367 & 0.037 \\
Speaker-instruction   & 1.000 & 0.367 & 0.033 \\
\midrule
Chance                & 0.08  & 0.17  & 0.011 \\
\bottomrule
\end{tabular}
\end{table}

\paragraph{Operator-A still extracts a real, bounded gain.}
The flatness of \emph{conditioning} does not mean the embedding is inert: choosing the readout itself---an Operator-A ($\cZA$) action over layer and pooling---yields a measurable emotion gain. \Cref{tab:opA} shows mean-pool emotion accuracy of $0.393$ at L0 (CI $[0.340,0.447]$) rising to $0.490$ under attentive-statistics pooling at L16 (CI $[0.433,0.550]$), a delta of $+0.097$ ($\approx +25\%$ relative). At a test size of $300$ the CIs touch, but the dev-selected seed-$42$ point ($0.513$) strengthens the effect. Speaker remains at chance ($0.047\!\to\!0.063$) and content stays at ceiling ($0.987\!\to\!0.997$), confirming that Operator-A reallocates expressivity it can reach but cannot manufacture speaker information the frozen space does not encode. On MInDS-14 en-US intent ($14$ intents, chance $0.071$), intent \emph{presence} is real---baseline $\approx 0.25$, CI $[0.195,0.304]$, far above chance---yet \emph{steerability} is again unsupported (selected conditioning emotion, delta $+0.016$, overlapping CIs, not diagonally dominant). Operator-A thus delivers a genuine but bounded gain, consistent with a near-flat single-model space.

\begin{table}[t]
\centering
\caption{In-house preliminary results. Operator-A ($\cZA$) layer/pooling selection on CREMA-D TEST (seed $42$, $1000$-bootstrap CIs). The emotion gain is real but bounded; speaker stays at chance, content at ceiling.}
\label{tab:opA}
\begin{tabular}{lccc}
\toprule
Factor & Mean-pool @L0 & Attentive-stats @L16 & $\Delta$ \\
\midrule
Emotion (CI)  & 0.393 \,[0.340,0.447] & 0.490 \,[0.433,0.550] & $+0.097$ \\
Speaker       & 0.047 & 0.063 & $+0.016$ \\
Content       & 0.987 & 0.997 & $+0.010$ \\
\bottomrule
\end{tabular}
\end{table}

\paragraph{Operator-B task-level selection already yields large gains.}
The contrast that motivates the agent action space is stark. Where single-model conditioning is flat, \emph{task-level} Operator-B ($\cZB$) selection---choosing among generative candidates under a verifiable reward, the training-free RL move of this series---produces large, well-separated improvements (\Cref{tab:opB}). On SER/emotion the embedding-side gain carries over ($0.393\!\to\!0.490$, $+0.097$). On SLURP intent, accuracy rises $0.550\!\to\!0.880$ ($+0.330$, CI $[0.288,0.374]$). On MInDS-14 intent, $0.883\!\to\!0.972$ ($+0.089$), with an independent raw-schema reproduction improving $0.852\!\to\!0.984$ ($+0.132$, CI $[0.082,0.187]$; $25$ fixes / $1$ regression, RTX 5090, $n{=}182$). On URO spoken-QA, $0.380\!\to\!0.715$ ($+0.335$, CI $[0.265,0.405]$), with a conservative reranking pass lifting this to $0.845$ at \emph{zero} regressions. These paired-bootstrap gains (seed $42$) are an order of magnitude larger than anything single-model conditioning offers, and the Lean-backed guarantees T1--T6 (proofs in \texttt{proofs/tfrl}) bound their reward spread and gain.

\begin{table}[t]
\centering
\caption{In-house preliminary results. Cross-team Operator-B ($\cZB$) sample-level training-free-RL gains (seed $42$, paired bootstrap). Task-level selection yields large, well-separated improvements, in contrast to the flat single-model conditioning of \Cref{tab:cremad-matrix}.}
\label{tab:opB}
\begin{tabular}{lccc}
\toprule
Task & Baseline & After Op-B & $\Delta$ (CI) \\
\midrule
Emotion / SER (CREMA-D)        & 0.393 & 0.490 & $+0.097$ \\
SLURP intent                   & 0.550 & 0.880 & $+0.330$ \,[0.288,0.374] \\
MInDS-14 intent                & 0.883 & 0.972 & $+0.089$ \\
\quad repro (raw$\to$policy)   & 0.852 & 0.984 & $+0.132$ \,[0.082,0.187] \\
URO spoken-QA                  & 0.380 & 0.715 & $+0.335$ \,[0.265,0.405] \\
\quad $+$ conservative rerank  & 0.715 & 0.845 & (0 regressions) \\
\bottomrule
\end{tabular}
\end{table}

\paragraph{Reading the evidence.}
Taken together, the three tables make a single quantitative argument. Within one frozen model, instruction conditioning is near-flat---an empirical instance of the OSA-1 regime, in which the reward spread $\spread$ available to a single-model search distribution $\qz$ is too small to support a useful gain $\gain$. Operator-A can still reach a real but bounded readout gain ($+0.097$ on emotion), and Operator-B task-level selection already delivers far larger gains ($+0.330$ on SLURP, $+0.132$ on MInDS-14, $+0.335$ on URO). The implication for design is direct: the productive action does not live in conditioning a single frozen model but in selecting and composing across operators and episodes---precisely the agent action space this work targets, with the verifiable-reward acceptance gate of \Cref{sec:feas:gate} as its control law. We report all intervals and label these numbers as preliminary; the CI overlaps on the smaller emotion effects are stated explicitly and temper the strength of the Operator-A claim while leaving the Operator-B contrast unambiguous.

% ===== section 08-system (writer wordcount 1180) =====
\section{Proposed System: A Self-Evolving Omni Speech Agent}
\label{sec:system}

We now instantiate the theory of \Cref{sec:osa} as a concrete, buildable agent. The design is governed by a single discipline: every component that \emph{evolves} at inference time does so only through a verifiable-reward acceptance gate operating inside a $\beta$-KL trust region, never through a gradient on the base model. This is what separates our agent from the prevailing self-evolving paradigm, whose append-only memories and self-generated skill libraries provably drift toward context collapse and error propagation \citep{ace2025,expfollowing2025,remembering2026}. We design \emph{against} that failure mode by realizing the optimization geometry of $\qs$ (the Gibbs optimum of \Cref{sec:osa}) as an external, mutable state whose update operator is contractive by construction.

\subsection{Architecture}
\label{subsec:arch}

The agent couples two \emph{frozen} omni models with one external evolving state. A frozen \emph{vector-omni} encoder \citep{omniembed2025,omniembedaudio2026} serves as the memory and skill index: it maps raw spoken turns into the embedding action space $\cZA$ of \Cref{sec:osa}, where retrieval is a tilting of the base index distribution $\qo$. A frozen \emph{generative omni} policy \citep{qwen3omni} serves as the response model, whose decoding occupies the generative action space $\cZB$. Neither model is fine-tuned; both are queried only. All adaptation lives in an external state $S$---a structured memory store and a skill library---which is mutated under verifiable speech rewards $\Rwd$ (WER/CER for ASR, exact-match for intent, speaker/emotion agreement for paralinguistics). Formally, the agent's inference-time degrees of freedom $\cZ = \cZA \times \cZB$ are conditioned on $S$, and learning is the trajectory $S_0 \to S_1 \to \cdots$ of \emph{state} updates rather than weight updates. Because $S$ is the only mutable object, the convergence question reduces entirely to the contractivity of the state-update operator, which we engineer directly rather than hope for.

\paragraph{Why two frozen models.} Separating the index ($\cZA$) from the policy ($\cZB$) lets each operator be tilted by an independent reward channel: memory writes are gated by retrieval-utility rewards, while response decoding is gated by task rewards. This factorization mirrors the Operator-A / Operator-B decomposition of \Cref{sec:osa} and lets us bound the two mutation rates separately (\Cref{subsec:control}).

\subsection{Memory}
\label{subsec:memory}

\paragraph{Three tiers over a bi-temporal graph.} Memory is a structured store with three tiers: \emph{episodic} (individual spoken turns), \emph{semantic} (per-speaker persona and stable facts), and \emph{procedural} (pointers into the skill library of \Cref{subsec:skills}). The tiers sit on a bi-temporal knowledge graph that records both event time and ingestion time \citep{arigraph2024,zep2025,amem2025}, so that retrieval can respect causality and the curator can reason about staleness. Episodic nodes consolidate into semantic nodes off the critical path.

\paragraph{Index on audio, not transcripts.} A central design choice is that the retrieval key is the frozen \emph{audio} embedding from the vector-omni encoder, not an ASR transcript \citep{wavrag2025,voxrag2025}. Transcribing before indexing discards exactly the paralinguistic structure (prosody, speaker identity, affect) that the agent must remember, and it injects ASR error into the key. Indexing on $\cZA$ keeps the key in the lossless acoustic space.

\paragraph{Dual-use composite keys.} Each memory entry carries a composite key
\[
k = (\text{speaker-ID},\ \text{emotion/SER},\ \text{turn-index}),
\]
extracted by frozen diarization, x-vector, and ECAPA-TDNN speaker embedding plus an SER head \citep{xvector2021,ecapatdnn2020}. Crucially these key fields are \emph{dual-use}: they are simultaneously the retrieval key and a source of verifiable reward. When the agent later answers a speaker- or emotion-conditioned query, agreement between the predicted attribute and the stored key field is a verifiable reward signal---no human label and no learned judge required. This closes the loop that \citet{memoryr1_2025} closes with a \emph{trained} memory manager: we obtain credit assignment from the keys themselves.

\paragraph{Retrieve--rerank.} Retrieval is two-stage. A fast first stage does late-interaction matching in $\cZA$ \citep{colbert2020}; a precise second stage reranks the shortlist by reward-weighted importance, weighting each candidate by its historical contribution to verifiable reward \citep{genagents2023}. The reranker is the discrete analogue of the importance weight $\exp(\Rwd/\beta)$ that defines $\qs$ in \Cref{sec:osa}.

\paragraph{Non-destructive curation with provenance.} Writes are mediated by a curator that emits one of $\{\textsc{Add}, \textsc{Update}, \textsc{Delete}, \textsc{Noop}\}$ \citep{mem0_2025}; \textsc{Noop} is first-class and dominant, which is what keeps the store from bloating. Salience decays over time on an Ebbinghaus schedule \citep{memorybank2023}, so unreinforced episodic traces fade unless consolidated. Every node carries provenance and a trust score; low-trust or anomalous writes are quarantined to resist memory poisoning \citep{agentpoison2024}, the dominant attack surface for evolving stores.

\paragraph{Gradient-free credit assignment.} The curator does \emph{not} learn its policy by gradient. Given a candidate set of memory-write operations, it admits the write only if a verifiable-reward best-of-$N$ comparison shows a positive reward delta against the no-write baseline. This replaces the trained manager of \citet{memoryr1_2025} with a training-free selection rule whose acceptance probability is exactly the tilting weight of \Cref{sec:osa}, inheriting the best-of-$N$ KL bounds of \citet{beirami2024bon,verdun2025softbon}.

\paragraph{$\beta$-KL trust region on mutation rate.} The first instantiation of the control law bounds the \emph{rate} of memory mutation: the curator may not move the store's distribution over keys by more than a $\beta$-controlled KL budget per episode,
\[
\KL{q_{S_{t+1}}}{q_{S_t}} \le \epsilon_{\mathrm{mem}}(\beta).
\]
This trust region is what formally prevents the context collapse documented for append-only agents \citep{ace2025,expfollowing2025}: a bounded-KL update operator is contractive in the sense required by \Cref{sec:osa}, so the state trajectory converges rather than drifting. Consolidation (episodic$\to$semantic compaction, decay sweeps, re-embedding) runs asynchronously off the critical path, in the spirit of sleep-time and asynchronous-retrieval memory services \citep{lettasleeptime,moshirag2026}, so the trust region never blocks live inference.

\subsection{Skills}
\label{subsec:skills}

\paragraph{The speech-skill object.} A skill is a structured \textsc{Skill.md}-style object $(C, \pi, T, R)$---a usage context $C$, a callable procedure $\pi$, a set of test cases $T$, and a verifiable reward $R$ \citep{sokskills2026,anthropicskills}. The procedure $\pi$ is drawn from a small set of speech-native bodies: generative error correction for ASR rare-word repair \citep{gerrarewords2025}, voice tool-use orchestration \citep{aura2025}, and paralinguistic-preserving speech-to-speech rewriting \citep{paras2s2025}. Each body is a callable over the frozen policy, never a fine-tune of it.

\paragraph{Verification gate.} A candidate skill is admitted to the library only if it produces a \emph{positive verifiable-reward delta} relative to the no-skill baseline on its test set $T$; a resident skill whose running delta turns negative is deprecated \citep{skillsbench2026}. This is the curation-quality result of the skills literature made operational: curated, verified skills dominate self-generated ones, and the gate is precisely what enforces curation. At selection time the agent retrieves skills body-first---matching the situation to skill \emph{bodies} rather than to their natural-language descriptions \citep{skillrouter2026}---which avoids the description-mismatch failure of name-based routing.

\paragraph{Gradient-free improvement under a Pareto trust region.} Skills improve without touching base weights through reflective, evolutionary search: GEPA-style Pareto editing of skill prompts and procedures \citep{gepa2025} and AFlow-style workflow search over skill compositions \citep{aflow2024}. The governing constraint is the \emph{second} instantiation of our control law: an edited skill is accepted only if it is \emph{Pareto non-regressive}---it must not lose verifiable reward on any task in its competency set. Pareto non-regression is a $\beta$-KL trust region on the skill library's behavior: the post-edit policy may not move outside a bounded-KL neighborhood of the pre-edit policy on any verified task, which guarantees monotone non-degradation and rules out the reflection-plateau and skill-thrash pathologies \citep{reflexion2023,lats2023}.

\paragraph{Bloat control and composition.} Library growth is contained by polymorphic abstraction---collapsing near-duplicate skills into parameterized families \citep{polyskill2025}---and by CRAFT-style validate-and-deduplicate at admission \citep{craft2023}. Genuinely new competence accretes by AWM-style snowballing, where verified skills compose into higher-order workflows that are themselves gated \citep{awm2024}. Thus the library grows in capability while its effective description length stays bounded, the skill-side analogue of the memory decay schedule.

\subsection{The Unified Control Law}
\label{subsec:control}

The two subsystems are governed by one law, stated once:
\begin{quote}
\emph{Admit an inference-time state mutation iff it yields a positive verifiable-reward delta against the no-mutation baseline (the acceptance gate), and only within a $\beta$-KL trust region that bounds how far the induced action distribution may move (the stability constraint).}
\end{quote}
The acceptance gate is the discrete realization of the reward tilting that defines the Gibbs optimum $\qs$, inheriting its $O(1/N)$ regret and KL--reward trade-off from \Cref{sec:osa}; the trust region is the regularizer $\beta\,\KL{\cdot}{\qo}$ that keeps the search distribution $\qz$ from collapsing. The law has exactly two instantiations: the memory-mutation-rate bound $\epsilon_{\mathrm{mem}}(\beta)$ of \Cref{subsec:memory}, and the skill-library Pareto non-regression bound of \Cref{subsec:skills}. Both are bounded-KL, hence both yield contractive state-update operators, and together they operationalize the credit-assigned convergence guarantee established in \Cref{sec:osa}: the external state trajectory $S_0 \to S_1 \to \cdots$ converges because each admitted step strictly improves verifiable reward while remaining inside a $\beta$-controlled neighborhood, with no base-model gradient anywhere in the loop.

% ===== section 09-plan (writer wordcount 1186) =====
\section{Research Plan}
\label{sec:plan}

We organize the empirical program as a partially-observed sequential decision problem over experiments: each phase is selected to maximize value-of-information per unit of compute, so that cheap, high-diagnostic measurements gate the expensive ones. The ordering is deliberate. We first close the loop on a single GPU with frozen weights (\Cref{sec:plan} Phase~0), because no downstream claim about training-free RL can be trusted until the verifiable-reward gate is demonstrably faithful end-to-end. We then build the missing measurement instrument (Phase~1), since the central scientific question---whether reward-guided, inference-time optimization over the action space $\cZ$ can \emph{compound} paralinguistic competence across sessions---cannot be answered without a benchmark that does not yet exist. Only then do we run the minimal self-evolving loop (Phase~2) and ablate its controls (Phase~3). Throughout, the only quantities that drive selection are verifiable rewards $\Rwd$, and the locked test set is touched exactly once.

\subsection{Phase 0: Closed-Loop Bring-Up}
\label{sub:phase0}

The objective of Phase~0 is a verified, reproducible inference-time loop on the in-place RTX~5090 (Blackwell, sm\_120; torch \texttt{cu128}) with \emph{no weight updates and no structural change}. The frozen stack couples an omni-embedding model (Omni-Embed-Nemotron, \citealp{omniembed2025,omniembedaudio2026}) for the embedding operator $\cZA$ with Qwen3-Omni \citep{qwen3omni} served via GGUF/\texttt{llama.cpp} for the generative operator $\cZB$. Concretely, this realizes the two-operator action space $\cZ = \cZA \cup \cZB$ used throughout: $\cZA$ selects pooling/probe read-outs of a frozen representation; $\cZB$ is reward-guided decoding (best-of-$N$, reranking, soft-BoN) over text.

We reproduce two known frozen baselines as correctness anchors: SER accuracy $\approx 0.49$ and intent accuracy $\approx 0.25$ under zero-shot prompting of the backbone. Matching these within bootstrap tolerance certifies that data loading, prompt templating (\texttt{speechrl\_common.models}), and decoding are wired correctly. We then verify the verifiable-reward gate: for ASR we instrument WER-based exact rewards, and for SER/SID we instrument label-exact rewards (\texttt{speechrl\_common.rl}), confirming that (i) best-of-$N$ selection under $\Rwd$ strictly dominates greedy on a held-in probe, consistent with the $O(1/N)$ regret/winrate scaling of soft and hard best-of-$N$ \citep{verdun2025softbon,beirami2024bon,aminian2025smoothing}, and (ii) the implied search distribution $\qz$ tilts the base $\qo$ toward the Gibbs optimum $\qs \propto \qo \exp(\Rwd/\beta)$ as $\beta$ is lowered, with the partition function $\Zpart$ and objective $\Fobj{\cdot}$ tracked as in the KL-as-Bayes view \citep{korbak2022kl}. The exit criterion is a green end-to-end run logging $\gain$ (reward gain over greedy) and $\spread$ (reward max$-$min) per utterance to local MLflow, plus a passing \texttt{pytest} smoke suite. Phase~0 buys the highest information per dollar: it eliminates instrumentation confounds before any benchmark or loop claim is made.

\subsection{Phase 1: A Speech Cross-Session Paralinguistic Memory Benchmark}
\label{sub:phase1}

\paragraph{The gap.} Existing long-horizon memory benchmarks are text-only (\citealp{longmemeval2024,locomo2024}; the agentic memory line \citealp{amber2026,anatomymem2026}), and audio memory work targets within-document or RAG retrieval \citep{wavrag2025,voxrag2025,audiomc2025}. No raw-audio, \emph{cross-session}, paralinguistically-keyed memory benchmark exists. This is the open contribution: a benchmark whose queries can only be answered by retaining \emph{who spoke} and \emph{how they sounded} across sessions, not by transcribing content.

\paragraph{Minimal spec.} (i) \emph{Stimuli}: multi-session spoken dialogues delivered as raw audio, with $\geq 3$ sessions per persona separated by simulated time gaps. (ii) \emph{Queries}: speaker-ID-keyed and SER-keyed (e.g., ``which speaker was angry in the previous call,'' ``has this caller's affect changed since session~1''), so the answer is a function of paralinguistic state, not lexical content---guarding against the read-not-listen shortcut \citep{voxparadox2026,listenortranscribe2025}. (iii) \emph{Ground truth}: verifiable SID/SER labels, enabling exact rewards rather than judge models. (iv) \emph{Probes}: a knowledge-update probe (a speaker's emotion/role changes across sessions, testing \textsc{update}/\textsc{delete} semantics \citealp{mem0_2025}), a contamination probe \citep{contaminationsurvey2024,taskcontamination2024}, and a memory-poisoning probe \citep{agentpoison2024,mempoison2026}. (v) \emph{Construction}: rather than collecting new audio, we re-key an existing diarized corpus (SLURP/SUPERB-style intent+speaker material, \citealp{slurp2020,superb2021,dynsuperb2023}) with \emph{frozen} diarization (pyannote, \citealp{pyannotecode,personalvad2019,xvector2021}) and a frozen SER labeler, then validate label fidelity on a human-checked subset. This keeps the benchmark cheap, reproducible, and license-clean.

\paragraph{Evaluation design.} Scoring uses verifiable rewards only, reported with paired bootstrap confidence intervals over utterances \citep{bisani2004bootstrap}, standard frozen splits \citep{gorman2019splits}, and an audited memory-eval rubric \citep{anatomymem2026}. We additionally report \emph{maintenance latency} (wall-clock cost of the memory \textsc{add}/\textsc{update} cycle) and run every metric across $\geq 2$ backbones (Qwen3-Omni and Qwen2.5-Omni, \citealp{qwen25omni}) to separate benchmark signal from backbone idiosyncrasy. Phase~1 is sequenced second because it is the precondition for any compounding claim, yet is independent of the loop's internals, so it can be built and frozen before Phase~2 risks overfitting to it.

\subsection{Phase 2: Minimal Self-Evolving Loop}
\label{sub:phase2}

Phase~2 instantiates the smallest loop that could plausibly compound: a paralinguistic \emph{memory} store (frozen-embedding keys with \textsc{add}/\textsc{update}/\textsc{delete} \citealp{mem0_2025,memoryos2025}), a \emph{skill gate} that admits a candidate skill only when a held-out verifiable reward improves \citep{skillsbench2026,coevoskills2026}, and a \emph{trust region} that bounds the per-episode change in the search distribution via $\KL{\qz}{\qo} \le \epsilon$, capping the inference-time tilt to prevent reward over-optimization \citep{gao2023overopt,skalse2022rewardhacking}. The loop is gradient-free: credit is assigned by reward comparison across rollouts \citep{jitrl2026,ahmadian2024rloo}, and selection over $\cZB$ uses reranking/MBR among $N$ samples \citep{ichihara2025mbr,mbrasr2025}.

We \emph{pre-register} success and kill criteria before running. Success: a positive, bootstrap-significant cross-episode slope on SER, intent, and spoken-QA reward, with the gate's admit-rate stable and $\KL{\qz}{\qo}$ inside the trust region. Kill: no slope after a fixed episode budget, or admit-rate collapse, or trust-region saturation (drift). The headline measurement is \emph{compounding}: reward as a function of episode index, contrasted against a single-shot baseline (best-of-$N$ with an empty, non-persistent memory), isolating the contribution of cross-session state from raw test-time compute \citep{snell2024ttscompute,zuo2025ttrl}.

\subsection{Phase 3: Scale and Ablate}
\label{sub:phase3}

Phase~3 stress-tests each control by removal, with pre-registered directional predictions. (i) \emph{Remove the skill gate}: we predict SkillsBench-style degradation as unvetted skills accumulate \citep{skillsbench2026,ace2025}. (ii) \emph{Remove the trust region}: we predict drift, i.e., $\KL{\qz}{\qo}$ growth correlated with falling locked-test reward, the inference-time analogue of reward hacking \citep{gao2023overopt,remembering2026}. (iii) \emph{Append-only vs.\ curated memory}: we predict append-only error propagation and context collapse \citep{expfollowing2025,ace2025}, with curated \textsc{update}/\textsc{delete} recovering it \citep{mem0_2025}. Each ablation is a single-variable contrast against the Phase~2 configuration, scaled across $N\in\{4,8,16,32\}$ and across backbones to chart the reward-vs-compute and reward-vs-tilt frontiers.

\subsection{Evaluation Protocol}
\label{sub:protocol}

The protocol is fixed in advance and applies uniformly. (1) \emph{Rewards}: verifiable only (WER/exact-match/label-exact); no LLM-judge in any selection or reporting path, avoiding self-preference and position bias \citep{selfpreference2024,positionbias2024}. (2) \emph{Test discipline}: the locked test set is evaluated exactly once at the end of each phase; all tuning and model selection occur on a disjoint dev split \citep{dwork2015reusable,roelofs2019overfitting}. (3) \emph{Demonstration hygiene}: in-context demonstrations are drawn disjoint from both dev and test, since demonstrations alter format more than accuracy \citep{min2022rethinking,alice2026}, and we report true few-shot results \citep{perez2021truefewshot}. (4) \emph{Uncertainty}: every comparison reports paired bootstrap CIs \citep{bisani2004bootstrap,dror2018hitchhiker} over $\geq 5$ seeds \citep{henderson2018drlmatters}, with a power analysis fixing the minimum detectable effect \citep{card2020power}. (5) \emph{Multiplicity}: family-wise correction across the ablation grid. (6) \emph{Negative results are reported}, and the full configuration follows the reproducibility checklist \citep{pineau2020reproducibility,dodge2019showyourwork}.

\subsection{Milestones, Compute, and Reproducibility}
\label{sub:milestones}

Compute is already provisioned: an in-place RTX~5090 with torch \texttt{cu128}, all five models (including the Qwen3-Omni GGUF and Omni-Embed-Nemotron \citealp{omniembedweights}) and all 28 datasets resident on WSL ext4 disk, frozen to a pinned manifest. Every run is reproducible by construction: pinned data revisions $+$ fixed seed $+$ a single \texttt{reproduce:} command logged to MLflow, so any reported number maps to one invocation.

\begin{table}[t]
\centering
\small
\begin{tabular}{llll}
\hline
\textbf{Milestone} & \textbf{Phase} & \textbf{Deliverable} & \textbf{Exit gate} \\
\hline
M0 & 0 & Closed loop on frozen $\cZA{+}\cZB$ & SER$\,{\approx}\,0.49$, intent$\,{\approx}\,0.25$ reproduced \\
M1 & 0 & Verifiable-reward gate verified & best-of-$N$ $>$ greedy, CI-significant \\
M2 & 1 & Cross-session paralinguistic benchmark & frozen splits $+$ human-checked labels \\
M3 & 1 & Probe suite (update/contam./poison) & probes pass face validity, 2 backbones \\
M4 & 2 & Minimal self-evolving loop & pre-registered success/kill adjudicated \\
M5 & 2 & Compounding vs.\ single-shot result & positive bootstrap slope or documented null \\
M6 & 3 & Control ablations & gate/TR/append-only contrasts complete \\
M7 & 3 & Frontier sweep ($N$, backbone) & reward-vs-compute/tilt curves $+$ release \\
\hline
\end{tabular}
\caption{Value-of-information--ordered milestones. Each gate is a verifiable, bootstrap-tested criterion; the locked test set is touched once per phase.}
\label{tab:milestones}
\end{table}

The release artifact bundles the benchmark builder, the loop configuration, pinned manifests, and per-milestone \texttt{reproduce:} commands, so that all of \Cref{tab:milestones} can be re-derived end-to-end on equivalent hardware. This ordering guarantees that the expensive compounding study (M4--M7) is only attempted after the instrument (M2--M3) and the gate (M0--M1) are independently certified, maximizing the information returned for the fixed single-GPU budget.

% ===== section 10-discussion (writer wordcount 1287) =====
\section{Threats to Validity and Limitations}
\label{sec:threats}

We hold this work to the standard of self-criticism it advocates, and enumerate the threats in
descending order of how much they could change the paper's claims. We distinguish \emph{threats to
the thesis} (does optimization-space adequacy actually hold?) from \emph{threats to the moat} (is the
contribution durable and novel?) and \emph{threats to the empirics} (do the measurements support the
conclusions?).

\paragraph{The moat is an absence of evidence (fast-follower risk).}
Our central feasibility argument rests on the observation that no prior work couples training-free,
reward-guided inference-time optimization with a frozen omni model's disentangled embeddings as the
search variable. An absence of evidence is structurally weaker than evidence of absence: it is
falsified the moment a competing group publishes the same coupling, and the individual ingredients
--- best-of-$N$ and reranking \citep{beirami2024bon,verdun2025softbon}, controlled decoding
\citep{mudgal2024controlled}, no-gradient credit assignment \citep{jitrl2026}, test-time RL
\citep{zuo2025ttrl,li2025tpo} --- are all public. We therefore do not claim a defensible technical
monopoly; we claim a \emph{first-mover synthesis} whose value is the formalization
(\Cref{sec:threats} notwithstanding, the OSA theorems) and the machine-checked guarantees, not the
mechanism itself. A reader should weight this contribution as integrative rather than as a primitive
discovery.

\paragraph{The mechanism is not novel; it is a domain transfer.}
The Gibbs-optimal tilting $\qs \propto \qo\exp(\Rwd/\beta)$ and its KL-regularized variational
characterization are textbook results in the alignment-as-Bayesian-inference literature
\citep{korbak2022kl,yang2024asymptotics}, and the best-of-$N$/KL frontier is now sharply understood
\citep{beirami2024bon,gui2024bonbon}. Our contribution is to instantiate this machinery over the
\emph{embedding} action space $\cZA$ of a frozen speech model rather than the token space, and to
argue that the relevant degrees of freedom $\cZ$ already carry the pretrained paralinguistic and
semantic structure \citep{frozenllmparaling2024,salmonn2023}. This is a transfer claim, and transfer
claims fail when the target domain violates a hidden precondition of the source domain --- precisely
the risks catalogued next.

\paragraph{OSA-2 idealizes context isolation as exact separability.}
\Cref{sec:threats} aside, OSA-2 models a multi-key agent reward as a sum of per-context-isolated
components and proves adequacy of the operator-A search distribution under \emph{exact} additive
separability of $\Rwd$ across the content, speaker, emotion, and language/intent factors. Real agent
rewards are at best \emph{approximately} separable: an ASR-conditioned downstream reward leaks into
the emotion channel through prosody, and intent classification is not independent of transcription
quality. Under a residual coupling of magnitude $\epsilon$ in the cross-factor Hessian, the
Gibbs optimum of the separable surrogate is displaced from the true optimum, and the gain $\gain$
we certify is an upper bound that degrades as $\Rwd$ departs from additivity. We have not
characterized this degradation quantitatively; the gap between idealized and approximate separability
is the single largest theory-to-practice risk, and the empirical separability of the embedding
factors \citep{xiao2020whatnot,wang2020alignuniform,leace2023} is an assumption our probes support
only locally, not globally.

\paragraph{OSA-3b proves the finite-separable optimum; finite-time convergence over the full space is open.}
OSA-3b establishes that, restricted to a finite, exactly separable action set, the reward-guided
update reaches the Gibbs optimum $\qs$ and the gate behaves as a stabilizing control law. The
operational convergence we actually run, however, is JitRL's \citep{jitrl2026} no-gradient credit
assignment, whose guarantee is \emph{asymptotic under slow drift} of the rollout distribution. The
two regimes do not compose into a finite-time statement: we do \emph{not} have a monotone-improvement
guarantee for the gated update over the full, continuous, only-approximately-separable space at a
finite horizon. This is an open problem, not an oversight, and it bears on the stability claims that
agentic deployments care about most \citep{ace2025,expfollowing2025,remembering2026}. Until it is
closed, the control-law framing of the gate is a design principle backed by an idealized proof, not a
deployment-grade certificate.

\paragraph{The isolated Hoeffding lemma and the T2 order-statistics \texttt{sorry} are load-bearing assumptions.}
Our machine-checked development contains two acknowledged gaps. The isolated Hoeffding lemma is
stated and used but its concentration constant is assumed rather than derived from the actual
reward's bounded-difference structure; if the per-context reward is heavier-tailed than the assumed
sub-Gaussian envelope, the sample-complexity constants are optimistic. The T2 order-statistics result
that underwrites the best-of-$N$ spread $\spread$ retains a \texttt{sorry}: the asymptotics of the
maximum order statistic \citep{yang2024asymptotics,beirami2024bon} are invoked as a hypothesis, not
proved in our formalization. We flag both explicitly so that no reader mistakes ``machine-checked''
for ``unconditionally proved''; the checked artifact is sound \emph{relative to} these two
assumptions.

\paragraph{Reward hacking and Goodhart on verifiable speech rewards.}
Inference-time optimization against a fixed reward is, by construction, an over-optimization engine:
pushing $N$ or lowering $\beta$ moves probability mass toward whatever the reward measures, including
its defects \citep{gao2023overopt,skalse2022rewardhacking}. Verifiable speech rewards (WER, exact
match, intent accuracy) are more robust than learned reward models, but they are not immune: WER
rewards literal transcription and can be gamed by hedging or copying, and ``spurious rewards'' can
induce gains that do not reflect the intended capability \citep{spuriousrewards2025,rlvrincentive2025}.
Because we change no weights, the failure is recoverable (lower $N$, raise $\beta$) rather than baked
into parameters --- but the gate must be tuned against held-out, contamination-controlled splits
\citep{contaminationsurvey2024,gorman2019splits}, or the certified $\gain$ measures the benchmark, not
the skill.

\paragraph{Paralinguistic keys are dual-use: a mis-estimated label both mis-routes and mis-rewards.}
The same emotion/speaker estimate that \emph{routes} a query (selecting the context and skill) also
\emph{rewards} the rollout that produced it. A mis-estimated SER label therefore enters two coupled
loops with the same sign of error: it sends the query to the wrong context \emph{and} reinforces the
embedding that yielded the wrong label, a positive-feedback pathway that idealized separability hides.
This is sharpened by evidence that omni models tend to ``read, not listen'' on paralinguistic cues
\citep{voxparadox2026,listenortranscribe2025} and that speaker identity resists in-context correction
\citep{spkverifllm2026}. The mitigation --- decoupling the routing estimator from the reward
estimator --- is future work, not a property of the current gate.

\paragraph{Memory poisoning via acoustic spoofing.}
When the optimized embedding is written to agent memory, the acoustic channel becomes an injection
surface: a spoofed speaker or emotion can plant a poisoned key that later retrieval amplifies
\citep{agentpoison2024,mempoison2026}. Unlike text poisoning, the attack needs no lexical trigger ---
it rides paralinguistic features the model already over-trusts. We provide no audit defense here;
memory hygiene and provenance \citep{anatomymem2026} are required before the embedding-as-memory-key
design is safe to deploy.

\paragraph{Small-sample confidence intervals.}
Several of our preliminary results are reported on modest evaluation sets, where bootstrap intervals
\citep{bisani2004bootstrap} are wide and statistical power is low \citep{card2020power,dror2018hitchhiker}.
Effects near the resolution of a few hundred utterances should be read as directional, not confirmed;
adequately powered, seed-averaged \citep{henderson2018drlmatters} replication on the frozen,
contamination-controlled split is a precondition for any quantitative headline.

\section{Conclusion}
\label{sec:conclusion}

We have argued, and where possible machine-checked, a single thesis: that the inference-time degrees
of freedom of a \emph{frozen} speech/omni multimodal LLM constitute an action space $\cZ$ that is
\emph{adequate} for the optimization we ask of it --- that reward-guided search over $\cZ$, with no
gradient and no weight or structure change, can activate pretrained capability rather than instill it.
We formalized this as optimization-space adequacy in three parts. OSA-1 identifies $\cZ$ and its two
operators, the embedding/Operator-A space $\cZA$ and the generative/Operator-B space $\cZB$, and
characterizes the Gibbs optimum $\qs \propto \qo\exp(\Rwd/\beta)$ as the target of any KL-regularized
search $\Fobj{q}=\Exp_q[\Rwd]-\beta\,\KL{q}{\qo}$. OSA-2 shows that, under context isolation, an
operator-A search distribution is adequate for the multi-key agentic reward. OSA-3 (with OSA-3b)
establishes the separable-finite optimum and the gate's behavior as a control law on the
exploration--exploitation frontier. We are deliberate that these are proved under the idealizations
and the two named assumptions of \Cref{sec:threats}: the contribution is a rigorous \emph{scaffold},
honestly bounded, not an unconditional theorem.

The feasibility case is correspondingly measured. The moat is open rather than owned, but the
\emph{components are ready}: verifiable speech rewards, best-of-$N$ and reranking, no-gradient credit
assignment \citep{jitrl2026}, and a frozen omni embedding model \citep{omniembed2025,qwen25omni} whose
factors our probes find locally separable. The gate reads not as a heuristic but as a control law with
an idealized optimality proof, and our preliminary operator-A results on emotion, speaker, and
language/intent factors are consistent with --- though not yet powered to confirm --- the predicted
gain $\gain$. The plan that follows is to close the two formalization gaps, to replace exact
separability with a quantified approximate-separability bound, to decouple the routing and reward
estimators against the dual-use feedback risk, and to replicate on the frozen,
contamination-controlled split with adequate power.

We close by reaffirming the discipline that makes the claim falsifiable: a survey-first posture that
treats our own moat as an absence of evidence to be defended, not assumed; a commitment to
\emph{verifiable} rewards over learned proxies, precisely to blunt the Goodhart pressure that
inference-time optimization invites; and a no-gradient, no-weight-change constraint that keeps every
failure recoverable at inference rather than frozen into parameters. The promise of activating, rather
than retraining, the latent competence of frozen speech models is, we believe, both real and
near-term --- and it is worth stating only to the precision that the proofs and the measurements can
currently bear.

% ===== section 11-appendix (writer wordcount 1520) =====
\appendix

\section{Lean 4 Formalization}
\label{app:lean}

The core of our training-free RL theory is machine-checked in \textbf{Lean~4 + mathlib} (toolchain and mathlib both pinned to \texttt{v4.31.0}), in the project \texttt{proofs/tfrl} (Lean namespace \texttt{TfrlProofs}). The full library builds with a \emph{single documented} \texttt{sorry}, isolated to the order-statistics step of the best-of-$N$ KL bound (\Cref{app:lean-bon}); every other theorem, including the entire \texttt{TfrlProofs.OptSpace} module that grounds the optimization-space adequacy (OSA) claims of the main text, is \textbf{\texttt{sorry}-free}. Throughout, $\cZ$ is a \texttt{Fintype} (the inference-time action space), $\qo:\cZ\to\bbR$ is a strictly positive reference distribution with $\sum_z \qo(z)=1$, $\Rwd:\cZ\to\bbR$ is a reward, and $\beta>0$ a temperature.

\paragraph{Tilting primitives and the Gibbs inequality.}
The exponentially tilted optimum $\qs$, its partition function $\Zpart$, and the KL-regularized objective $\Fobj{\cdot}$ are defined exactly as in the main text. We reproduce the Lean signatures verbatim.

\begin{verbatim}
noncomputable def Zpart (q0 R : Z → ℝ) (β : ℝ) : ℝ := ∑ w, q0 w * Real.exp (R w / β)
noncomputable def qstar (q0 R : Z → ℝ) (β : ℝ) (z : Z) : ℝ :=
  q0 z * Real.exp (R z / β) / Zpart q0 R β
noncomputable def F (q0 R : Z → ℝ) (β : ℝ) (q : Z → ℝ) : ℝ :=
  (∑ z, q z * R z) - β * ∑ z, q z * Real.log (q z / q0 z)

theorem kl_nonneg {p r : Z → ℝ} (hp : ∀ z, 0 ≤ p z) (hr : ∀ z, 0 < r z)
    (hpsum : ∑ z, p z = 1) (hrsum : ∑ z, r z = 1) :
    0 ≤ ∑ z, p z * Real.log (p z / r z)
\end{verbatim}

\noindent
\texttt{kl\_nonneg} is Gibbs' inequality: the pointwise estimate $p(z)-r(z)\le p(z)\log(p(z)/r(z))$ (trivial when $p(z)=0$; otherwise $\log x \le x-1$ with $x=r(z)/p(z)$ via \texttt{Real.log\_le\_sub\_one\_of\_pos}) summed against $\sum p=\sum r=1$.

\paragraph{Master identity and T1.}
The decomposition $\Fobj{\qs}-\Fobj{q}=\beta\,\KL{q}{\qs}$ is the engine of optimality.

\begin{verbatim}
theorem F_sub_eq_beta_mul_kl [Nonempty Z] (hq0 : ∀ z, 0 < q0 z) (hβ : 0 < β)
    {q : Z → ℝ} (hq : ∀ z, 0 ≤ q z) (hqsum : ∑ z, q z = 1) :
    F q0 R β (qstar q0 R β) - F q0 R β q = β * ∑ z, q z * Real.log (q z / qstar q0 R β z)

theorem tilting_optimal [Nonempty Z] (hq0 : ∀ z, 0 < q0 z) (hβ : 0 < β)
    {q : Z → ℝ} (hq : ∀ z, 0 ≤ q z) (hqsum : ∑ z, q z = 1) :
    F q0 R β q ≤ F q0 R β (qstar q0 R β)
\end{verbatim}

\noindent
\textbf{T1} (\texttt{tilting\_optimal}) is then immediate: $\Fobj{\qs}-\Fobj{q}=\beta\,\KL{q}{\qs}\ge 0$ by \texttt{kl\_nonneg} and $\beta>0$.

\paragraph{T3 --- flat reward, no tilt.}
\begin{verbatim}
theorem qstar_eq_q0_of_const [Nonempty Z] {q0 : Z → ℝ} {β c : ℝ}
    (hq0sum : ∑ z, q0 z = 1) {R : Z → ℝ} (hR : ∀ z, R z = c) :
    ∀ z, qstar q0 R β z = q0 z
\end{verbatim}
\noindent
When $\Rwd(z)\equiv c$, $\Zpart=\exp(c/\beta)$ collapses and $\qs(z)=\qo(z)\exp(c/\beta)/\exp(c/\beta)=\qo(z)$: a degenerate action space with zero reward spread admits no steering.

\paragraph{T2 --- best-of-$N$ KL bound.}
\label{app:lean-bon}
\begin{verbatim}
noncomputable def klBoundBoN (N : ℕ) : ℝ := Real.log N - ((N : ℝ) - 1) / N
theorem klBoundBoN_nonneg {N : ℕ} (hN : 1 ≤ N) : 0 ≤ klBoundBoN N
theorem kl_best_of_n_le (klBoN : ℕ → ℝ) {N : ℕ}
    (hBeirami : klBoN N ≤ klBoundBoN N) :
    klBoN N ≤ Real.log N - ((N : ℝ) - 1) / N := hBeirami
\end{verbatim}
\noindent
We prove rigorously that $\spread_{\mathrm{BoN}}(N):=\log N-(N-1)/N\ge 0$ (it equals $\log N - 1 + 1/N$, nonnegative since $\log N\ge 1-1/N$). The order-statistics derivation $\KL{\pi_{\mathrm{BoN}}}{\qo}\le \texttt{klBoundBoN}\,N$ \citep{beirami2024bon} is the lone documented \texttt{sorry} (\texttt{klBoN\_le\_klBoundBoN\_TODO}); \texttt{kl\_best\_of\_n\_le} takes that estimate as a hypothesis and concludes.

\paragraph{T6 --- $O(\sqrt{\log N})$ regret.}
\begin{verbatim}
theorem regret_O_sqrt_log {gain R_max kl : ℝ} {N : ℕ} (hN : 1 ≤ N)
    (hRmax : 0 ≤ R_max) (hPinsker : gain ≤ R_max * Real.sqrt (kl / 2))
    (hKL : kl ≤ klBoundBoN N) :
    gain ≤ R_max * Real.sqrt (Real.log N / 2)
\end{verbatim}
\noindent
Chaining the Pinsker estimate $\gain\le \Rwd_{\max}\sqrt{kl/2}$ (hypothesis) with $kl\le\texttt{klBoundBoN}\,N\le\log N$ and monotonicity of $\sqrt{\cdot}$ yields $\gain = O(\sqrt{\log N})$.

\paragraph{The OSA suite (\texttt{TfrlProofs.OptSpace}, \texttt{sorry}-free).}
We define the optimization gain $\gain:=\Fobj{\qs}-\Fobj{\qo}$ and prove its structural properties.

\begin{verbatim}
noncomputable def gain (q0 R : Z → ℝ) (β : ℝ) : ℝ :=
  F q0 R β (qstar q0 R β) - F q0 R β q0

theorem gain_eq [Nonempty Z] (hq0 : ∀ z, 0 < q0 z) (hβ : 0 < β)
    (hq0sum : ∑ z, q0 z = 1) :
    gain q0 R β = β * Real.log (Zpart q0 R β) - ∑ z, q0 z * R z
theorem gain_nonneg [Nonempty Z] (hq0 : ∀ z, 0 < q0 z) (hβ : 0 < β)
    (hq0sum : ∑ z, q0 z = 1) : 0 ≤ gain q0 R β
theorem flat_no_gain [Nonempty Z] (hq0 : ∀ z, 0 < q0 z) (hβ : 0 < β)
    (hq0sum : ∑ z, q0 z = 1) {c : ℝ} (hR : ∀ z, R z = c) : gain q0 R β = 0
theorem gain_le_of_hoeffding [Nonempty Z] (hq0 : ∀ z, 0 < q0 z) (hβ : 0 < β)
    (hq0sum : ∑ z, q0 z = 1) {S : ℝ}
    (hHoeff : β * Real.log (Zpart q0 R β) - ∑ z, q0 z * R z ≤ S) :
    gain q0 R β ≤ S

theorem Zpart_product :
    Zpart (fun p : Z1 × Z2 => q01 p.1 * q02 p.2) (fun p => R1 p.1 + R2 p.2) β
      = Zpart q01 R1 β * Zpart q02 R2 β
theorem gain_product [Nonempty Z1] [Nonempty Z2] (hq01 : ∀ z, 0 < q01 z)
    (hq02 : ∀ z, 0 < q02 z) (hβ : 0 < β)
    (hq01sum : ∑ z, q01 z = 1) (hq02sum : ∑ z, q02 z = 1) :
    gain (fun p : Z1 × Z2 => q01 p.1 * q02 p.2) (fun p => R1 p.1 + R2 p.2) β
      = gain q01 R1 β + gain q02 R2 β
theorem qstar_product (p : Z1 × Z2) :
    qstar (fun p : Z1 × Z2 => q01 p.1 * q02 p.2) (fun p => R1 p.1 + R2 p.2) β p
      = qstar q01 R1 β p.1 * qstar q02 R2 β p.2

theorem rollout_deficit [Nonempty Z] (hq0 : ∀ z, 0 < q0 z) (hβ : 0 < β)
    {q : Z → ℝ} (hq : ∀ z, 0 ≤ q z) (hqsum : ∑ z, q z = 1) :
    F q0 R β q = F q0 R β (qstar q0 R β)
        - β * ∑ z, q z * Real.log (q z / qstar q0 R β z)
      ∧ 0 ≤ β * ∑ z, q z * Real.log (q z / qstar q0 R β z)
\end{verbatim}

\noindent
\texttt{gain\_eq} rewrites the gain as $\beta\log\Exp_{\qo}[e^{\Rwd/\beta}]-\Exp_{\qo}[\Rwd]$ (\textbf{OSA-1a}); \texttt{gain\_nonneg} (\textbf{OSA-1b}) is Jensen / Gibbs; \texttt{flat\_no\_gain} (\textbf{OSA-1c}) recovers T3 --- a single-model output-context (degenerate $\cZ$) affords no RL improvement; \texttt{gain\_le\_of\_hoeffding} (\textbf{OSA-1d}) isolates the analytic bound $S$, so $\spread\to 0\Rightarrow\gain\to 0$. The product lemmas (\textbf{OSA-2}, \textbf{OSA-3b}) show that over a context-isolated product $\cZ_1\times\cZ_2$ with $\qo=\qo^{(1)}\otimes\qo^{(2)}$ and separable $\Rwd(z_1,z_2)=\Rwd_1(z_1)+\Rwd_2(z_2)$, the gain is \emph{additive} and $\qs$ factorizes; iterating over $k$ isolated agents gives $\gain=\sum_i\gain_i$, growing with $k$. Finally \texttt{rollout\_deficit} (\textbf{OSA-3a}) shows any rollout policy $q$ sits below $\qs$ by exactly $\beta\,\KL{q}{\qs}\ge 0$: a naive rollout that fails to reach $\qs$ incurs a strictly positive deficit, while the $\beta$-KL trust region enforces the slow-drift precondition of credit-assigned tilting \citep{jitrl2026}.

\section{Extended Proofs}
\label{app:extended}

We give in full the two analytic results that the Lean development isolates as hypotheses or pure algebra.

\subsection{Hoeffding instance: the gain ceiling \texorpdfstring{$\spread^2/(8\beta)$}{spread2/8beta}}
\label{app:hoeffding}

\begin{lemma}[Bounded-reward gain ceiling]
\label{lem:hoeffding}
Let $\Rwd:\cZ\to[\Rwd_{\min},\Rwd_{\max}]$ with reward spread $\spread:=\Rwd_{\max}-\Rwd_{\min}$, and let $\qo$ be any distribution on $\cZ$ with $\Exp_{\qo}$ the corresponding expectation. Then for every $\beta>0$,
\[
\beta\,\log\Exp_{\qo}\!\big[e^{\Rwd/\beta}\big]-\Exp_{\qo}[\Rwd]\;\le\;\frac{\spread^2}{8\beta},
\]
and hence, by \texttt{gain\_le\_of\_hoeffding} with $S=\spread^2/(8\beta)$, $\;\gain(\qo,\Rwd,\beta)\le \spread^2/(8\beta)$.
\end{lemma}

\begin{proof}
Let $X:=\Rwd(z)/\beta$ for $z\sim\qo$, a random variable supported in $[a,b]$ with $a=\Rwd_{\min}/\beta$, $b=\Rwd_{\max}/\beta$, so $b-a=\spread/\beta$. Define the cumulant generating function $\psi(\lambda):=\log\Exp[e^{\lambda(X-\Exp X)}]$. We use Hoeffding's lemma: for a random variable $X$ supported in $[a,b]$ and any $\lambda\in\bbR$,
\[
\psi(\lambda)\;=\;\log\Exp\!\big[e^{\lambda(X-\Exp X)}\big]\;\le\;\frac{\lambda^2(b-a)^2}{8}.
\]
To prove this standard bound, set $\varphi(\lambda):=\log\Exp[e^{\lambda X}]$. Then $\varphi(0)=0$, $\varphi'(0)=\Exp X$, and $\varphi''(\lambda)=\operatorname{Var}_{\lambda}(X)$, where $\operatorname{Var}_{\lambda}$ is the variance under the tilted law $d\mu_\lambda\propto e^{\lambda x}d\mu$. Since $X\in[a,b]$ almost surely, the same holds $\mu_\lambda$-a.s., and a random variable bounded in an interval of length $b-a$ has variance at most $((b-a)/2)^2=(b-a)^2/4$ (the extremal case is the two-point distribution at the endpoints). Hence $\varphi''(\lambda)\le (b-a)^2/4$ for all $\lambda$. By Taylor's theorem with integral remainder, for some $\xi\in[0,\lambda]$,
\[
\varphi(\lambda)=\varphi(0)+\lambda\varphi'(0)+\tfrac12\lambda^2\varphi''(\xi)\le \lambda\,\Exp X+\frac{\lambda^2 (b-a)^2}{8},
\]
so $\psi(\lambda)=\varphi(\lambda)-\lambda\Exp X\le \lambda^2(b-a)^2/8$, which is Hoeffding's lemma.

Now take $\lambda=1$. Then
\[
\log\Exp[e^{X}]-\Exp[X]=\psi(1)\le \frac{(b-a)^2}{8}=\frac{1}{8}\Big(\frac{\spread}{\beta}\Big)^2=\frac{\spread^2}{8\beta^2}.
\]
Multiplying by $\beta>0$ and substituting $X=\Rwd/\beta$ (so $\Exp[X]=\Exp_{\qo}[\Rwd]/\beta$ and $\Exp[e^X]=\Exp_{\qo}[e^{\Rwd/\beta}]$) gives
\[
\beta\log\Exp_{\qo}[e^{\Rwd/\beta}]-\Exp_{\qo}[\Rwd]\le \frac{\spread^2}{8\beta}.
\]
By \texttt{gain\_eq}, the left side is exactly $\gain$, completing the proof.
\end{proof}

\Cref{lem:hoeffding} makes the scaling explicit: the inference-time RL gain is quadratically throttled by the reward spread, $\gain=O(\spread^2/\beta)$, so a flat or near-flat reward landscape (T3, \textbf{OSA-1c}) yields vanishing headroom regardless of search budget --- the analytic counterpart of reward over-optimization saturation \citep{gao2023overopt}.

\subsection{Product factorization (OSA-2) in full}
\label{app:product}

We verify the algebra underlying \texttt{Zpart\_product}, \texttt{meanR\_product}, and \texttt{gain\_product}. Let $\cZ=\cZ_1\times\cZ_2$, $\qo(z_1,z_2)=\qo^{(1)}(z_1)\,\qo^{(2)}(z_2)$ with $\sum\qo^{(1)}=\sum\qo^{(2)}=1$, and $\Rwd(z_1,z_2)=\Rwd_1(z_1)+\Rwd_2(z_2)$ (context isolation).

\begin{proposition}[Partition factorization]
$\Zpart(\qo,\Rwd,\beta)=\Zpart(\qo^{(1)},\Rwd_1,\beta)\cdot\Zpart(\qo^{(2)},\Rwd_2,\beta)$.
\end{proposition}
\begin{proof}
Expanding over the product index set and using $e^{(\Rwd_1(z_1)+\Rwd_2(z_2))/\beta}=e^{\Rwd_1(z_1)/\beta}e^{\Rwd_2(z_2)/\beta}$,
\[
\Zpart=\sum_{z_1,z_2}\qo^{(1)}(z_1)\qo^{(2)}(z_2)\,e^{\Rwd_1(z_1)/\beta}e^{\Rwd_2(z_2)/\beta}
=\Big(\sum_{z_1}\qo^{(1)}(z_1)e^{\Rwd_1(z_1)/\beta}\Big)\Big(\sum_{z_2}\qo^{(2)}(z_2)e^{\Rwd_2(z_2)/\beta}\Big),
\]
the last step by distributing the double sum (\texttt{Finset.sum\_mul\_sum}), which is the product of the two partition functions.
\end{proof}

\begin{proposition}[Mean reward additivity]
$\Exp_{\qo}[\Rwd]=\Exp_{\qo^{(1)}}[\Rwd_1]+\Exp_{\qo^{(2)}}[\Rwd_2]$.
\end{proposition}
\begin{proof}
For fixed $z_1$, $\sum_{z_2}\qo^{(1)}(z_1)\qo^{(2)}(z_2)(\Rwd_1(z_1)+\Rwd_2(z_2))=\qo^{(1)}(z_1)\Rwd_1(z_1)\sum_{z_2}\qo^{(2)}(z_2)+\qo^{(1)}(z_1)\sum_{z_2}\qo^{(2)}(z_2)\Rwd_2(z_2)=\qo^{(1)}(z_1)\Rwd_1(z_1)+\qo^{(1)}(z_1)\Exp_{\qo^{(2)}}[\Rwd_2]$, using $\sum\qo^{(2)}=1$. Summing over $z_1$ and again using $\sum\qo^{(1)}=1$ gives $\Exp_{\qo^{(1)}}[\Rwd_1]+\Exp_{\qo^{(2)}}[\Rwd_2]$.
\end{proof}

\begin{corollary}[Additive gain, OSA-2]
$\gain(\qo,\Rwd,\beta)=\gain(\qo^{(1)},\Rwd_1,\beta)+\gain(\qo^{(2)},\Rwd_2,\beta)$.
\end{corollary}
\begin{proof}
By \texttt{gain\_eq}, $\gain=\beta\log\Zpart-\Exp_{\qo}[\Rwd]$. Substituting the two propositions and $\log(xy)=\log x+\log y$ (legitimate since each partition function is strictly positive, \texttt{Zpart\_pos}),
\[
\gain=\beta\big(\log\Zpart_1+\log\Zpart_2\big)-\big(\Exp_{\qo^{(1)}}[\Rwd_1]+\Exp_{\qo^{(2)}}[\Rwd_2]\big)=\gain_1+\gain_2.
\]
Iterating over $k$ context-isolated components yields $\gain=\sum_{i=1}^k\gain_i$, strictly increasing in $k$ whenever each component has nonzero reward spread (\Cref{lem:hoeffding}). The companion identity $\qs(z_1,z_2)={\qs}^{(1)}(z_1)\,{\qs}^{(2)}(z_2)$ (\texttt{qstar\_product}, \textbf{OSA-3b}) follows by the same exponential split, certifying that per-component credit-assigned tilting reaches the global T1-optimum.
\end{proof}

\section{Notation}
\label{app:notation}

\begin{center}
\begin{tabular}{ll}
\hline
Symbol & Meaning \\
\hline
$\cZ$ & inference-time action space (degrees of freedom changing no weights) \\
$\cZA,\cZB$ & embedding/Operator-A space; generative/Operator-B space \\
$\qo$ & base reference policy $q_0$ (strictly positive, $\sum_z\qo(z)=1$) \\
$\qz$ & a search distribution $q$ over $\cZ$ \\
$\qs$ & Gibbs / tilted optimum $\qs(z)=\qo(z)e^{\Rwd(z)/\beta}/\Zpart$ \\
$\Rwd$ & reward function $\cZ\to\bbR$ \\
$\beta$ & temperature ($>0$); KL-regularization strength \\
$\Zpart$ & partition function $\sum_w \qo(w)e^{\Rwd(w)/\beta}$ \\
$\Fobj{q}$ & KL-regularized objective $\Exp_q[\Rwd]-\beta\,\KL{q}{\qo}$ \\
$\KL{p}{q}$ & Kullback--Leibler divergence \\
$\Exp$ & expectation $\mathbb{E}$ \\
$\gain$ & optimization gain $\Fobj{\qs}-\Fobj{\qo}\ge 0$ \\
$\spread$ & reward spread $\Rwd_{\max}-\Rwd_{\min}$ \\
$N$ & number of candidates in best-of-$N$ / search budget \\
\hline
\end{tabular}
\end{center}

\section{Verifiable Reward Functions}
\label{app:rewards}

The reward callables used as verifiable signals $\Rwd$ are exposed by the shared library \texttt{speechrl\_common.rl}, and are deliberately deterministic and contamination-resistant \citep{tulu3}. Heavy dependencies (\texttt{jiwer}, \texttt{scikit-learn}) are imported lazily inside each function so the package and its smoke tests import without the full stack.

\paragraph{Content / ASR (\texttt{rl.reward}).}
\texttt{wer(reference, hypothesis)} returns the word error rate after normalization; \texttt{asr\_reward} returns $\max(0,\,1-\mathrm{WER})\in[0,1]$, a bounded reward whose spread directly governs the gain ceiling of \Cref{lem:hoeffding}; \texttt{exact\_match\_reward} returns $\mathbf{1}[\text{normalized strings equal}]\in\{0,1\}$ for closed-form tasks (intent, keyword), matching the binary RLVR signal of T4-style plurality selection.

\paragraph{Retrieval / embedding (\texttt{rl.embedding\_metrics}).}
For the flagship embedding action space $\cZA$, \texttt{recall\_at\_k(query, gallery, labels\_q, labels\_g, k)} and \texttt{mean\_reciprocal\_rank(\dots)} score a frozen model's pooled embeddings against a labelled gallery under cosine similarity (\texttt{cosine\_sim\_matrix} with $\ell_2$ normalization); \texttt{retrieval\_reward(\dots, k)} packages recall@$k$ as a $[0,1]$ reward for speaker-ID / language retrieval.

\paragraph{Probe accuracy (\texttt{rl.probe}).}
\texttt{linear\_probe\_accuracy} / \texttt{knn\_probe\_accuracy} fit a frozen-feature classifier (held-out split) to score disentanglement of emotion (SER) and language/intent factors; \texttt{probe\_reward(\dots, kind)} returns the held-out accuracy as the reward, providing the verifiable signal for the content/speaker/emotion/language axes the flagship disentangles. Auxiliary metrics (\texttt{macro\_f1}, \texttt{bleu}, \texttt{chrf}, \texttt{eer}) in \texttt{rl.metrics} are used for evaluation only, not as RL rewards, to avoid coupling the reward and the test statistic \citep{skalse2022rewardhacking}.


\bibliography{references}

\end{document}
